%%%%%%%%%%%%%%%%%%%%%%%%%%%%%%%%%%%%%%%%%%%%%%%%%%%%%%%%%%%%%%%%%%%%%%%%%%%%%%%
%
% APPENDIX BBB: METHODOLOGICAL FOUNDATIONS OF PARAMETER ESTIMATION
% 
% The Epistemics of Numbers: From Empirical Data to Structured Priors
%
% EBF Framework – Behavioral Change Model
% FehrAdvice & Partners AG / BEATRIX Research Group
%
% Status: METHODOLOGY CORE (Required)
% Version: 1.0
% Category: METHOD
% Language: English
% Last Updated: 2026-01-06
%
%%%%%%%%%%%%%%%%%%%%%%%%%%%%%%%%%%%%%%%%%%%%%%%%%%%%%%%%%%%%%%%%%%%%%%%%%%%%%%%

\documentclass[11pt,a4paper]{article}

\usepackage{amsmath,amssymb,amsthm}
\usepackage{booktabs}
\usepackage{longtable}
\usepackage{graphicx}
\usepackage{hyperref}
\usepackage{natbib}
\usepackage{tcolorbox}
\usepackage{tikz}
\usepackage{array}

\theoremstyle{definition}
\newtheorem{definition}{Definition}[section]
\newtheorem{axiom}{Axiom}[section]
\newtheorem{theorem}{Theorem}[section]
\newtheorem{proposition}{Proposition}[section]
\newtheorem{example}{Example}[section]
\newtheorem{protocol}{Protocol}[section]

\title{Appendix EST: Methodological Foundations of Parameter Estimation\\
\large The Epistemics of Numbers: From Empirical Data to Structured Priors}
\author{EBF Framework\\BEATRIX Research Group\\FehrAdvice \& Partners AG}
\date{Version 1.0 -- January 2026}

\begin{document}

\maketitle

\begin{tcolorbox}[colback=blue!5!white, colframe=blue!75!black, title=Methodology Appendix]
\textbf{CORE-METHOD} | Answers: ``Where do the numbers come from?''

\begin{tabular}{ll}
\textbf{Status:} & CORE (Required) \\
\textbf{Dependencies:} & Appendix HOW (\ref{app:complementarity-how}) (Complementarity Theory) \\
 & Appendix LLM1 (\ref{app:llm-monte-carlo}) (LLM Monte Carlo) \\
 & Appendix OPS (\ref{app:operationalization}) (Empirical Operationalization) \\
\textbf{Key Output:} & The $\Theta$-Vector (Complete Parameter Set) \\
\textbf{Integration:} & Defines the ``Epistemic Status'' of all model inputs \\
\end{tabular}
\end{tcolorbox}

\tableofcontents
\newpage

%==============================================================================
% PART I: EPISTEMOLOGICAL FOUNDATIONS
%==============================================================================


%##############################################################################
%
%  APPENDIX BBB: THE EPISTEMICS OF NUMBERS
%  =======================================
%  From Empirical Data to Structured Priors
%
%  CORE Question: WHERE do the numbers come from?
%
%##############################################################################

%==============================================================================
% PART 1: FRONT MATTER
%==============================================================================

%------------------------------------------------------------------------------
% 1.1 HEADER BLOCK (Metadata)
%------------------------------------------------------------------------------

\begin{tcolorbox}[colback=red!5!white, colframe=red!75!black, 
    title=\textbf{Appendix EST: The Epistemics of Numbers}]

\begin{tabular}{@{}ll@{}}
\textbf{Full Title:} & The Epistemics of Numbers: From Empirical Data to Structured Priors \\
\textbf{Category:} & CORE (Fundamental Theory) \\
\textbf{Scope:} & A4 (Extended: 50--100 KB) \\
\textbf{Current Size:} & 80 KB \\
\textbf{CORE Question:} & \textbf{WHERE} do the numbers come from? \\
\textbf{Symbol:} & $E(\theta), \hat{\theta}, \sigma_\theta$ \\
\textbf{Version:} & 1.0 \\
\textbf{Last Updated:} & January 2026 \\
\textbf{Status:} & Complete \\
\end{tabular}

\tcblower

\textbf{Key Contribution:} This appendix provides the complete methodological 
framework for parameter estimation in EBF, resolving the fundamental 
trilemma between theoretical richness, data scarcity, and quantification need 
through a three-tier hybrid strategy.

\end{tcolorbox}

%------------------------------------------------------------------------------
% 1.2 CHAPTER LINKAGE BOX
%------------------------------------------------------------------------------

\begin{tcolorbox}[colback=blue!5!white, colframe=blue!75!black, 
    title=\textbf{Chapter Linkages}]

\textbf{Primary Chapter Links:}

\begin{center}
\begin{tabular}{@{}clp{7cm}@{}}
\toprule
\textbf{Chapter} & \textbf{Title} & \textbf{How BBB Supports} \\
\midrule
Ch. 3 & Model Implementation & Provides all parameter values \\
Ch. 4 & Calibration & Methodological foundation \\
Ch. 5 & Interventions & Effect size estimates \\
Ch. 7 & Complementarity & $\Gamma$-matrix coefficients \\
Ch. 8 & Context & $\delta$-tensor for context modulation \\
\bottomrule
\end{tabular}
\end{center}

\textbf{How to Use:}
\begin{itemize}[nosep]
    \item When implementing the model (Ch. 3), look up parameter values in \S\ref{sec:bbb-dictionary}
    \item When calibrating (Ch. 4), follow the Tier Integration Protocol in \S\ref{sec:bbb-three-tier}
    \item When estimating effects (Ch. 5), use the confidence scores from \S\ref{sec:bbb-uncertainty}
\end{itemize}

\end{tcolorbox}

%------------------------------------------------------------------------------
% 1.3 CROSS-REFERENCE MAP
%------------------------------------------------------------------------------

\begin{tcolorbox}[colback=gray!5!white, colframe=gray!75!black, 
    title=\textbf{Cross-Reference Map}]

\begin{center}
\small
\begin{tabular}{@{}c|ccccccc@{}}
\toprule
& \textbf{B} & \textbf{V} & \textbf{C} & \textbf{AAA} & \textbf{AN} & \textbf{E} & \textbf{AU} \\
\midrule
\textbf{BBB $\leftarrow$} & Theory & $\Psi$-def & FEPSDE & Levels & LLM-MC & Data & Axioms \\
\textbf{BBB $\rightarrow$} & $\Gamma,\Lambda,\delta$ & -- & $w_d$ & Params & Valid. & -- & -- \\
\midrule
\textbf{Link} & \cellcolor{red!20}Strong & \cellcolor{red!20}Strong & \cellcolor{yellow!20}Med & \cellcolor{yellow!20}Med & \cellcolor{red!20}Strong & \cellcolor{yellow!20}Med & \cellcolor{green!20}Weak \\
\bottomrule
\end{tabular}
\end{center}

\textbf{Legend:} 
\colorbox{red!20}{Strong} = Bidirectional dependency \quad
\colorbox{yellow!20}{Med} = One-way dependency \quad
\colorbox{green!20}{Weak} = Reference only

\end{tcolorbox}

%==============================================================================
% FORMAL COMPLIANCE BLOCK
%==============================================================================

% --- Terminology Reference ---
\begin{tcolorbox}[colback=gray!5!white, colframe=gray!50!black]
\textbf{Terminology:} See Appendix GLS (Glossary) for standard definitions.
Key terms in this appendix: Epistemic status tags (EMP/THR/LLM/ILL/HYP),
$\Theta$-vector (complete parameter set), prior specification,
LLM Monte Carlo, calibration trilemma, confidence scores.
\end{tcolorbox}

% --- Epistemic Status Tags ---
\begin{tcolorbox}[colback=blue!5!white, colframe=blue!75!black,
    title=\textbf{Epistemic Status: The Tag Hierarchy}]

\renewcommand{\arraystretch}{1.3}
\begin{tabular}{@{}llccc@{}}
\toprule
\textbf{Tag} & \textbf{Description} & \textbf{Evidence Level} & \textbf{Weight $w_s$} & \textbf{Stars} \\
\midrule
\textbf{[EMP]} & Empirically validated & Peer-reviewed experiments & 1.0 & $\star\star\star\star\star$ \\
\textbf{[THR]} & Theoretically derived & Formal model prediction & 0.8 & $\star\star\star\star$ \\
\textbf{[LLM]} & LLM Monte Carlo & Structured estimation & 0.6 & $\star\star\star$ \\
\textbf{[ILL]} & Illustrative & Example/demonstration & 0.3 & $\star\star$ \\
\textbf{[HYP]} & Hypothetical & Speculative/assumed & 0.1 & $\star$ \\
\addlinespace
\multicolumn{5}{@{}l}{\textit{Framework statistics:}} \\
Total parameters & In $\Theta$-vector & 163+ & & \\
EMP coverage & Parameters with [EMP] & $\approx 40\%$ & & \\
LLM coverage & Parameters with [LLM] & $\approx 35\%$ & & \\
\bottomrule
\end{tabular}

\vspace{0.5em}
\textbf{The Calibration Trilemma:} Precision vs. Coverage vs. Verifiability ---
choose two. EBF prioritizes Coverage + Transparency over false Precision.

\end{tcolorbox}

% --- Bibliography Reference ---
\noindent\textbf{Full citations:} See \texttt{bibliography/bcm\_master.bib} \\
\textbf{Key sources:} Gelman et al. (2013) Bayesian methods; DellaVigna (2009) meta-analysis; OpenAI (2023) LLM capabilities

\vspace{1em}

%------------------------------------------------------------------------------
% 1.4 ABSTRACT
%------------------------------------------------------------------------------

\begin{tcolorbox}[colback=white, colframe=black,
    title=\textbf{Abstract}]

EBF requires 185 parameters, yet direct empirical estimates exist for fewer 
than 20\%. This appendix presents a \textbf{three-tier hybrid estimation strategy} 
that resolves this challenge: (1) Hard empirical calibration from experimental 
data and meta-analyses, (2) Structured LLM Monte Carlo using large language 
models as ``compressed scientific knowledge,'' and (3) Theoretical constraints 
from axiomatic foundations. We detail five estimation methods, provide the 
complete parameter dictionary with epistemic status tags, and demonstrate the 
full workflow through a worked example estimating $\gamma_{FP}$ for Switzerland. 
The result is a transparent, updatable parameter infrastructure where every 
number has a documented provenance.

\textit{(98 words)}

\end{tcolorbox}

%------------------------------------------------------------------------------
% 1.5 QUICK REFERENCE BOX
%------------------------------------------------------------------------------

\begin{tcolorbox}[colback=green!5!white, colframe=green!75!black, 
    title=\textbf{Quick Reference}]

\textbf{Core Concepts:}
\begin{itemize}[nosep]
    \item Three Tiers: Empirical ($E \approx 0.9$) $>$ LLM-MC ($E \approx 0.8$) $>$ Axiomatic ($E = 1.0$ for constraints)
    \item Five Methods: A (Econometric), B (LLM-MC), C (Experimental), D (Meta-Analytic), E (Milgrom)
    \item Parameter Universe: 185 total (15 $\Gamma$ + 28 $\Lambda$ + 120 $\delta$ + 22 other)
\end{itemize}

\textbf{Key Equations:}
\begin{align}
E(\theta) &= \frac{\sum_{s} w_s \cdot q_s \cdot r_s}{\sum_{s} w_s} \tag{Epistemic Status} \\
\gamma_{dd'}(\boldsymbol{\Psi}) &= \bar{\gamma}_{dd'} + \sum_{k} \delta_{dd'k} \cdot (\Psi_k - 0.5) \tag{Context Modulation} \\
\text{Prov}(\theta) &= \langle \hat{\theta}, \sigma_\theta, E(\theta), \mathcal{S}(\theta), t(\theta), \mathcal{M}(\theta) \rangle \tag{Provenance Tuple}
\end{align}

\textbf{Jump to:}
\begin{itemize}[nosep]
    \item Parameter values $\rightarrow$ \S\ref{sec:bbb-dictionary}
    \item Estimation methods $\rightarrow$ \S\ref{sec:bbb-methods}
    \item Worked example $\rightarrow$ \S\ref{sec:bbb-example}
    \item Axiom system $\rightarrow$ \S\ref{sec:bbb-axiom-system}
    \item \textbf{Complete methodology navigation $\rightarrow$ Appendix PM (REF-PARAMETH)}
\end{itemize}

\end{tcolorbox}




%==============================================================================
% CRITICAL DISCLAIMER - EPISTEMIC STATUS OF PARAMETER VALUES
%==============================================================================

\begin{tcolorbox}[colback=red!10!white, colframe=red!75!black, 
    title=\textbf{⚠ CRITICAL DISCLAIMER: Epistemic Status of Parameter Values}]

\textbf{IMPORTANT:} The parameter values presented in this appendix fall into 
four distinct categories with very different epistemic statuses:

\begin{center}
\renewcommand{\arraystretch}{1.4}
\begin{tabular}{@{}clcp{6cm}@{}}
\toprule
\textbf{Tag} & \textbf{Category} & \textbf{Count} & \textbf{Meaning} \\
\midrule
\cellcolor{green!20}\texttt{EMP} & Empirical & ~15 & Direct estimates from published studies \\
\cellcolor{blue!20}\texttt{LLM} & LLM-MC & ~25 & Synthetic priors from LLM Monte Carlo \\
\cellcolor{yellow!20}\texttt{ILL} & Illustrative & ~100 & Plausible values for demonstration \\
\cellcolor{gray!20}\texttt{THR} & Theoretical & ~45 & Derived from axiomatic constraints \\
\bottomrule
\end{tabular}
\end{center}

\tcblower

\textbf{What this means for users:}

\begin{itemize}[nosep]
\item \texttt{EMP} values can be used with confidence (always check original source)
\item \texttt{LLM} values are reasonable starting points requiring local validation
\item \texttt{ILL} values are \textbf{placeholders} showing the structure---\textbf{DO NOT use in production without replacement by real estimates}
\item \texttt{THR} values are boundary conditions and constraints, not point estimates
\end{itemize}

\textbf{The purpose of illustrative values:} They demonstrate \textit{how} the 
framework works and \textit{what kind of} parameters are needed. They are 
scaffolding for future empirical work, not final answers.

\end{tcolorbox}



\label{app:bbb-estimation}
% -----------------------------------------------------------------------------
% APPENDIX SCOPE BOX
% -----------------------------------------------------------------------------

\begin{tcolorbox}[
    colback=orange!5!white,
    colframe=orange!75!black,
    title={\textbf{Appendix Scope: Ziel / In-Scope / Out-of-Scope / Constraints}},
    fonttitle=\bfseries
]

\textbf{Ziel (Objective):}
\begin{itemize}[nosep]
    \item [TODO: Define objective for Unknown Title]
\end{itemize}

\textbf{In-Scope:}
\begin{itemize}[nosep]
    \item The Three-Tier Hybrid Estimation Strategy
    \item The Five Methods of Parameter Estimation
    \item The Complete Parameter Dictionary
    \item Bayesian Updating: From Priors to Posteriors
    \item Handling Uncertainty and Heterogeneity
\end{itemize}

\textbf{Out-of-Scope (delegiert an andere Appendices):}
\begin{itemize}[nosep]
    \item \ref{app:complementarity-how} $\rightarrow$ Appendix HOW
    \item \ref{app:operationalization} $\rightarrow$ Appendix OPS
    \item Glossary $\rightarrow$ Appendix GLS
    \item \ref{app:complementarity-how} $\rightarrow$ Appendix HOW
    \item \ref{app:complementarity-how} $\rightarrow$ Appendix HOW
\end{itemize}

\textbf{Constraints (Anwendungsgrenzen):}
\begin{itemize}[nosep]
    \item [TODO: Define when this appendix does NOT apply]
    \item [TODO: Define required prerequisites]
    \item [TODO: Define known limitations]
\end{itemize}

\textbf{Lieferobjekte (Deliverables):}
\begin{itemize}[nosep]
    \item 8 Table(s)
    \item 59 Equation(s)
    \item Worked Example(s)
\end{itemize}

\end{tcolorbox}


\section{Epistemological Foundations}
\label{sec:bbb-epistemology}

The EBF framework makes quantitative predictions about behavioral phenomena 
that have traditionally been treated as qualitative. This requires a rigorous 
epistemological stance on parameter estimation---what we call the 
\textbf{Epistemics of Numbers}.

\subsection{The Fundamental Challenge}

Structural behavioral models face a trilemma:

\begin{enumerate}
    \item \textbf{Theoretical Richness:} We know behavior is context-dependent 
    ($U = f(\mathbf{d}, \boldsymbol{\Psi})$), with complementarities between 
    utility dimensions and context factors.
    
    \item \textbf{Data Scarcity:} We rarely have simultaneous measurements of all 
    8 $\Psi$-dimensions, all 6 FEPSDE utilities, and behavioral outcomes in a 
    single dataset.
    
    \item \textbf{Quantification Need:} To be actionable, the model must output 
    numbers (effect sizes, intervention priorities), not just directional 
    metaphors (``context matters'').
\end{enumerate}

This trilemma has led to two unsatisfying responses in the literature:
\begin{itemize}
    \item \textbf{Over-simplification:} Ignoring context (``one effect size fits all'')
    \item \textbf{Under-specification:} Acknowledging heterogeneity without 
    quantifying it (``effects vary by context'')
\end{itemize}

EBF offers a third path: \textbf{structured estimation with explicit 
epistemic status}.

\subsection{The Solution: A Hierarchy of Evidence}

We resolve the trilemma by rejecting the binary distinction between ``measured'' 
and ``unknown''. Instead, we adopt a continuous \textbf{Hierarchy of Evidence}, 
where every parameter $\theta$ carries an \textbf{Epistemic Tag} $E(\theta) \in [0, 1]$ 
indicating its source reliability.

\begin{definition}[Epistemic Tag]
\label{def:epistemic-tag}
For any parameter $\theta$ in the EBF model:
\begin{equation}
E(\theta) = \frac{\text{Evidence Quality} \times \text{Relevance}}{\text{Uncertainty}}
\end{equation}
where:
\begin{itemize}
    \item $E(\theta) \approx 1.0$: Direct experimental measurement
    \item $E(\theta) \in [0.7, 0.9]$: Meta-analytic or LLM-synthesized estimate
    \item $E(\theta) \in [0.3, 0.7]$: Calibrated from aggregate data
    \item $E(\theta) < 0.3$: Theoretical constraint or informed guess
\end{itemize}
\end{definition}

\subsection{The Parameter Universe}

EBF requires estimation of the following parameter classes:

\begin{table}[h]
\centering
\begin{tabular}{llcc}
\toprule
\textbf{Parameter Class} & \textbf{Symbol} & \textbf{Count} & \textbf{Defined In} \\
\midrule
Utility-Utility Complementarity & $\Gamma = [\gamma_{dd'}]$ & 15 & Appendix HOW (\ref{app:complementarity-how}), Teil II \\
Context-Context Complementarity & $\Lambda = [\lambda_{kl}]$ & 28 & Appendix HOW (\ref{app:complementarity-how}), Teil IV \\
Context Modulation & $\delta_{dd'k}$ & 120 & Appendix HOW (\ref{app:complementarity-how}), Teil III \\
Base Utility Weights & $w_d$ & 6 & Appendix WAT (\ref{app:fepsde-matrix}) \\
Reference Points & $r_d$ & 6 & Appendix WAT (\ref{app:fepsde-matrix}) \\
Discount Factor & $\delta$ & 1 & Chapter 3 \\
Tipping Thresholds & $\Psi^*_k$ & 8 & Appendix HOW (\ref{app:complementarity-how}), Teil VI \\
Coherence Threshold & $K^*$ & 1 & Appendix HOW (\ref{app:complementarity-how}), Teil I \\
\midrule
\textbf{Total} & $\Theta$ & \textbf{185} & \\
\bottomrule
\end{tabular}
\caption{The complete parameter universe of EBF}
\label{tab:parameter-universe}
\end{table}

%==============================================================================
% PART II: THE THREE-TIER HYBRID STRATEGY
%==============================================================================

\section{The Three-Tier Hybrid Estimation Strategy}
\label{sec:bbb-three-tier}

Parameters in EBF are derived from three distinct source tiers, 
prioritized lexicographically. When a parameter can be estimated from 
multiple tiers, we use the highest-reliability source available.

\subsection{Tier 1: Hard Empirical Calibration (Gold Standard)}

\begin{tcolorbox}[colback=green!5!white, colframe=green!75!black]
\textbf{Reliability:} $E(\theta) \approx 0.9 - 1.0$ \\
\textbf{Method:} Direct econometric estimation from experimental or high-quality survey data \\
\textbf{Reference:} Appendix OPS (\ref{app:operationalization}) (Operationalization), Chapter 4
\end{tcolorbox}

\textbf{Sources:}
\begin{itemize}
    \item Meta-analyses of experimental economics (ultimatum games, dictator games, 
    public goods experiments)
    \item Large-scale surveys (World Values Survey, European Social Survey, 
    Gallup World Poll)
    \item Randomized controlled trials with behavioral outcomes
    \item Administrative data with behavioral markers
\end{itemize}

\textbf{Parameters estimated via Tier 1:}
\begin{itemize}
    \item Base discount factor $\delta$ (from intertemporal choice experiments)
    \item Social preference parameters (from dictator/ultimatum games)
    \item Risk aversion coefficients (from lottery experiments)
    \item Key $\gamma$ coefficients with direct evidence (e.g., $\gamma_{FP}$ from 
    health economics)
\end{itemize}

\begin{example}[Health-Wealth Complementarity from PSID]
Using Panel Study of Income Dynamics:
\begin{align}
\text{Wellbeing}_{it} &= \alpha + \beta_F \text{Income}_{it} + \beta_P \text{Health}_{it} \nonumber \\
&\quad + \gamma_{FP} (\text{Income} \times \text{Health})_{it} + \mu_i + \varepsilon_{it}
\end{align}
Estimate: $\hat{\gamma}_{FP} = 0.23$ (SE = 0.04), $E(\gamma_{FP}) = 0.92$
\end{example}

\subsection{Tier 2: Structured LLM Monte Carlo (Synthetic Priors)}

\begin{tcolorbox}[colback=yellow!5!white, colframe=yellow!75!black]
\textbf{Reliability:} $E(\theta) \approx 0.7 - 0.9$ \\
\textbf{Method:} LLM-MC Protocol with perspective rotation and cross-validation \\
\textbf{Reference:} Appendix LLM1 (\ref{app:llm-monte-carlo}) (LLM Monte Carlo Methodology)
\end{tcolorbox}

Large Language Models serve as ``compressed knowledge engines'' that have 
absorbed vast amounts of scientific literature, case studies, and empirical 
patterns. We use them not as oracles but as \textit{structured prior generators}.

\textbf{The LLM-MC Protocol} (detailed in Appendix LLM1 (\ref{app:llm-monte-carlo})):
\begin{enumerate}
    \item \textbf{Persona Specification:} Define behavioral agent with context 
    $\boldsymbol{\Psi}$ and current utility profile
    \item \textbf{Scenario Variation:} Present choices that reveal 
    complementarity/substitutability
    \item \textbf{Perspective Rotation:} Query from multiple ``expert'' viewpoints 
    (economist, psychologist, practitioner)
    \item \textbf{Monte Carlo Sampling:} Generate $N = 1000$ estimates with 
    varied temperatures
    \item \textbf{Cross-Validation:} Compare against available empirical benchmarks
\end{enumerate}

\textbf{Parameters estimated via Tier 2:}
\begin{itemize}
    \item Most $\gamma_{dd'}$ coefficients without direct experimental evidence
    \item Context modulation coefficients $\delta_{dd'k}$
    \item The full $\Lambda$-matrix of context-context interactions
\end{itemize}

\begin{table}[h]
\centering
\small
\begin{tabular}{lcccc}
\toprule
\textbf{Parameter} & \textbf{LLM-MC} & \textbf{Empirical} & \textbf{Correlation} & $E(\theta)$ \\
\midrule
$\gamma_{FP}$ & 0.21 & 0.23 & 0.91 & 0.85 \\
$\gamma_{ES}$ & 0.34 & 0.31 & 0.88 & 0.82 \\
$\gamma_{FE}$ & 0.08 & 0.11 & 0.82 & 0.78 \\
$\gamma_{PD}$ & 0.15 & 0.18 & 0.85 & 0.80 \\
\bottomrule
\end{tabular}
\caption{Validation of Tier 2 estimates against Tier 1 benchmarks}
\label{tab:tier2-validation}
\end{table}

\subsection{Tier 3: Theoretical and Axiomatic Constraints}

\begin{tcolorbox}[colback=gray!5!white, colframe=gray!75!black]
\textbf{Reliability:} $E(\theta) = 1.0$ (by definition), but low granularity \\
\textbf{Method:} Logical deduction from axioms \\
\textbf{Reference:} Appendix AWA (\ref{app:axiom-formalization}) (Axiom Formalization)
\end{tcolorbox}

Some parameters are determined by logical necessity rather than empirical 
estimation:

\textbf{Examples:}
\begin{itemize}
    \item \textbf{Time budget constraint:} Activities competing for time exhibit 
    strict substitutability: $\gamma_{\text{Time}} = -1$
    \item \textbf{Identity matrix:} $\gamma_{dd} = 0$ (a dimension cannot be 
    complementary to itself in the interaction sense)
    \item \textbf{Symmetry:} $\gamma_{dd'} = \gamma_{d'd}$ (by Young's theorem)
    \item \textbf{Monotonicity bounds:} $\gamma_{dd'} \in [-1, 1]$ (normalized)
\end{itemize}

\subsection{The Tier Integration Protocol}

\begin{protocol}[Tier Integration]
\label{prot:tier-integration}
For each parameter $\theta$:
\begin{enumerate}
    \item \textbf{Check Tier 1:} Is there direct experimental evidence?
    \begin{itemize}
        \item If yes: Use Tier 1 estimate, set $E(\theta) > 0.9$
    \end{itemize}
    \item \textbf{Check Tier 3:} Is the parameter logically constrained?
    \begin{itemize}
        \item If yes: Use axiomatic value, set $E(\theta) = 1.0$ for the constraint
    \end{itemize}
    \item \textbf{Default to Tier 2:} Use LLM-MC estimate
    \begin{itemize}
        \item Set $E(\theta)$ based on LLM consensus and cross-validation
    \end{itemize}
    \item \textbf{Document:} Record source tier and uncertainty in Parameter Dictionary
\end{enumerate}
\end{protocol}

%==============================================================================
% PART III: THE FIVE ESTIMATION METHODS
%==============================================================================


%==============================================================================
% JUSTIFICATION FOR TIER THRESHOLDS
%==============================================================================

\subsection{Justification of Tier Thresholds}
\label{sec:bbb-tier-thresholds}

\textbf{Why 0.9 and 0.7?}

The tier thresholds are \textbf{not arbitrary}---they derive from decision-theoretic 
considerations about acceptable error rates in policy applications.

\begin{tcolorbox}[colback=gray!5!white, colframe=gray!75!black, 
    title=\textbf{Derivation of Tier Thresholds}]

\textbf{Tier 1 Threshold ($E \geq 0.9$):}

For high-stakes interventions (e.g., national policy), we require:
\begin{itemize}[nosep]
\item False positive rate $\alpha < 0.05$ (standard scientific threshold)
\item Power $(1-\beta) > 0.8$ (Cohen's convention)
\item Combined: $E \approx (1-\alpha)(1-\beta) = 0.95 \times 0.80 \approx 0.76$
\end{itemize}

We round \textbf{up} to $E \geq 0.9$ because:
\begin{enumerate}[nosep]
\item Policy errors are costly and hard to reverse
\item Replication crisis suggests published estimates are often inflated
\item Conservative threshold protects against overconfidence
\end{enumerate}

\textbf{Tier 2 Threshold ($E \geq 0.7$):}

For medium-stakes applications (pilots, localized interventions):
\begin{itemize}[nosep]
\item Acceptable error rate $\alpha < 0.10$
\item Minimum power $(1-\beta) > 0.7$
\item Combined: $E \approx 0.90 \times 0.70 \approx 0.63$
\end{itemize}

We round \textbf{up} to $E \geq 0.7$ to maintain meaningful separation from Tier 3.

\textbf{Tier 3 ($E < 0.7$):}

Parameters below this threshold are exploratory. They may inform research 
directions but should not drive resource allocation decisions.

\end{tcolorbox}

\textbf{Sensitivity Analysis:}

The exact thresholds (0.9, 0.7) can be adjusted based on:
\begin{itemize}[nosep]
\item \textbf{Stakes}: Higher stakes $\Rightarrow$ higher thresholds
\item \textbf{Reversibility}: Reversible interventions tolerate lower $E$
\item \textbf{Domain norms}: Some fields accept more uncertainty
\end{itemize}

\begin{center}
\small
\begin{tabular}{@{}lccc@{}}
\toprule
\textbf{Application Context} & \textbf{Tier 1} & \textbf{Tier 2} & \textbf{Tier 3} \\
\midrule
National policy & $\geq 0.95$ & $\geq 0.80$ & $< 0.80$ \\
Standard (EBF default) & $\geq 0.90$ & $\geq 0.70$ & $< 0.70$ \\
Exploratory research & $\geq 0.80$ & $\geq 0.60$ & $< 0.60$ \\
\bottomrule
\end{tabular}
\end{center}



\section{The Five Methods of Parameter Estimation}
\label{sec:bbb-methods}

Within the three-tier framework, we employ five distinct estimation methods 
depending on data availability and parameter type.

\subsection{Method A: Econometric Estimation from Panel Data}

\textbf{Use case:} When rich longitudinal data exists with variation in both 
$X$ and outcomes $Y$.

\subsubsection{The Interaction Regression Approach}

For a panel of individuals $i$ over time $t$:

\begin{equation}
Y_{it} = \alpha + \sum_d \beta_d X_{d,it} + \sum_{d < d'} \gamma_{dd'} X_{d,it} X_{d',it} + \mu_i + \tau_t + \varepsilon_{it}
\label{eq:interaction-regression}
\end{equation}

\begin{proposition}[Identification Conditions]
\label{prop:gamma-id}
$\gamma_{dd'}$ is identified if:
\begin{enumerate}
    \item Sufficient variation in both $X_d$ and $X_{d'}$
    \item Interaction $X_d \cdot X_{d'}$ not collinear with main effects
    \item Selection controlled via fixed effects or IV
\end{enumerate}
\end{proposition}

\subsubsection{Context-Dependent Complementarity}

To estimate $\gamma_{dd'}(\boldsymbol{\Psi})$:

\begin{equation}
Y_{it} = \alpha + \sum_d \beta_d X_{d,it} + \sum_{d < d'} \left( \bar{\gamma}_{dd'} + \sum_k \delta_{dd'k} \Psi_{k,it} \right) X_{d,it} X_{d',it} + \varepsilon_{it}
\end{equation}

This three-way interaction model requires substantial sample sizes 
($N > 50$ per cell).

\subsection{Method B: LLM Monte Carlo Estimation}

\textbf{Use case:} When direct data is scarce but scientific knowledge exists 
in distributed form across the literature.

\subsubsection{The LLM-MC Protocol}

\begin{definition}[LLM Monte Carlo for $\Gamma$-Estimation]
\begin{enumerate}
    \item \textbf{Persona specification:} Define agent with context $\boldsymbol{\Psi}$ 
    and utility profile $(u_F, u_E, u_P, u_S, u_D, u_E)$
    \item \textbf{Scenario variation:} Present choices trading off dimensions 
    $d$ vs. $d'$ at varying levels
    \item \textbf{Choice elicitation:} Record simulated choices across $N$ scenarios
    \item \textbf{Parameter recovery:} Estimate $\gamma_{dd'}$ via maximum likelihood
\end{enumerate}
\end{definition}

\textbf{Scenario Template:}
\begin{verbatim}
"[Persona] faces a choice between:
 Option A: +Δ in {d}, -ε in {d'}
 Option B: -ε in {d}, +Δ in {d'}
 Option C: +δ in both {d} and {d'}
 
 Current levels: {d} = X_d, {d'} = X_{d'}
 Context: Ψ = (Ψ_I, ..., Ψ_F)"
\end{verbatim}

Under utility $U = u_d + u_{d'} + \gamma_{dd'} u_d u_{d'}$, the choice probability 
for Option C reveals $\gamma_{dd'}$.

\begin{axiom}[LLM Calibration Validity]
\label{ax:bbb-llm}
\textbf{(BBB-LLM)} LLM Monte Carlo estimates are valid when:
\begin{enumerate}
    \item Persona specification matches target population
    \item Scenario framing avoids obvious biases
    \item Cross-validation with empirical data shows $r > 0.8$
\end{enumerate}
\end{axiom}

\subsection{Method C: Experimental Estimation (Factorial Design)}

\textbf{Use case:} When randomized intervention is feasible.

\begin{definition}[2$\times$2 Factorial Design for Complementarity]
Randomly assign to:
\begin{center}
\begin{tabular}{c|cc}
& $X_{d'} = 0$ & $X_{d'} = 1$ \\
\hline
$X_d = 0$ & Control & Treatment $d'$ only \\
$X_d = 1$ & Treatment $d$ only & Both treatments \\
\end{tabular}
\end{center}
\end{definition}

\begin{theorem}[Complementarity from Factorial Design]
\begin{equation}
\gamma_{dd'} = \underbrace{[Y_{11} - Y_{10}]}_{\text{Effect of } d' | d=1} - \underbrace{[Y_{01} - Y_{00}]}_{\text{Effect of } d' | d=0}
\end{equation}
$\gamma_{dd'} > 0$ indicates complementarity.
\end{theorem}

\subsection{Method D: Meta-Analytic Bayesian Calibration}

\textbf{Use case:} When multiple studies exist but with heterogeneous designs.

Start with prior $\gamma_{dd'} \sim N(\mu_0, \sigma_0^2)$ based on theory.
Update with each study $s$:

\begin{equation}
p(\gamma | \text{data}_1, \ldots, \text{data}_S) \propto p(\gamma) \prod_{s=1}^{S} p(\text{data}_s | \gamma)
\end{equation}

\begin{table}[h]
\centering
\small
\begin{tabular}{lccccc}
\toprule
\textbf{Parameter} & \textbf{Prior $\mu_0$} & \textbf{Prior $\sigma$} & \textbf{Post. $\mu$} & \textbf{Post. $\sigma$} & \textbf{Studies} \\
\midrule
$\gamma_{FP}$ & 0.20 & 0.10 & 0.22 & 0.03 & 47 \\
$\gamma_{ES}$ & 0.30 & 0.15 & 0.32 & 0.04 & 31 \\
$\gamma_{FE}$ & 0.05 & 0.10 & 0.09 & 0.02 & 28 \\
\bottomrule
\end{tabular}
\caption{Bayesian meta-analytic estimates}
\label{tab:bayesian-meta}
\end{table}

\subsection{Method E: Milgrom Diagnostic Tests}

\textbf{Use case:} When we need to \textit{detect} complementarity without 
directly estimating $\gamma$.

The ten theorems from \citet{brynjolfsson2013complementarity} provide indirect 
tests:

\begin{table}[h]
\centering
\small
\begin{tabular}{p{2.5cm}p{4cm}p{5cm}}
\toprule
\textbf{Theorem} & \textbf{Prediction} & \textbf{Test} \\
\midrule
1. Pairwise & All pairs supermodular & $\binom{n}{2}$ interaction tests \\
2. Clustering & Bimodal practices & Hartigan's dip test \\
3. Coherent Search & Same-direction changes & Sign test on $\Delta x$ \\
4. Monotone CS & $\theta \uparrow \Rightarrow$ all $x_i \uparrow$ & IV regression \\
7. Le Chatelier & Long-run $>$ short-run & Compare $\beta^{LR}$ vs $\beta^{SR}$ \\
8. Momentum & All-or-nothing change & Hazard models \\
\bottomrule
\end{tabular}
\caption{Milgrom diagnostic tests}
\label{tab:milgrom}
\end{table}

\begin{definition}[Complementarity Diagnostic Index]
\begin{equation}
\text{CDI} = \frac{1}{6} \sum_{t=1}^{6} \mathbb{1}_{\text{Test } t \text{ passes}}
\end{equation}
CDI $\geq 0.83$ indicates strong complementarity evidence.
\end{definition}

\begin{axiom}[Milgrom Diagnostic Validity]
\label{ax:bbb-mil}
\textbf{(BBB-MIL)} Milgrom tests are valid when:
\begin{enumerate}
    \item $N > 30$ per cell
    \item Exogenous variation for MCS test
    \item Panel data for Le Chatelier/Momentum
\end{enumerate}
\end{axiom}

\subsection{Method Selection Guide}

\begin{axiom}[Estimation Method Selection]
\label{ax:bbb-est}
\textbf{(BBB-EST)} Choose method based on:
\begin{enumerate}
    \item \textbf{Data availability:} Experimental $>$ Panel $>$ Cross-section $>$ LLM
    \item \textbf{Causal identification:} Experimental $>$ Panel+FE $>$ LLM $>$ Meta
    \item \textbf{Cost efficiency:} LLM $>$ Meta $>$ Panel $>$ Experimental
    \item \textbf{Context specificity:} LLM $>$ Experimental $>$ Panel $>$ Meta
\end{enumerate}
\end{axiom}

%==============================================================================
% PART IV: THE PARAMETER DICTIONARY
%==============================================================================

\section{The Complete Parameter Dictionary}
\label{sec:bbb-dictionary}

This section provides the definitive reference for all EBF parameters.

\subsection{$\Gamma$-Matrix: Utility-Utility Complementarity}

\begin{longtable}{llcccc}
\toprule
\textbf{Pair} & \textbf{Name} & \textbf{$\hat{\gamma}$} & \textbf{SE} & \textbf{Tier} & $E(\theta)$ \\
\midrule
\endfirsthead
\toprule
\textbf{Pair} & \textbf{Name} & \textbf{$\hat{\gamma}$} & \textbf{SE} & \textbf{Tier} & $E(\theta)$ \\
\midrule
\endhead
$\gamma_{FE}$ & Financial-Emotional & 0.09 & 0.02 & 1+2 & 0.88 \\
$\gamma_{FP}$ & Financial-Physical & 0.23 & 0.04 & 1 & 0.92 \\
$\gamma_{FS}$ & Financial-Social & 0.11 & 0.03 & 2 & 0.82 \\
$\gamma_{FD}$ & Financial-Digital & 0.14 & 0.03 & 2 & 0.79 \\
$\gamma_{FEc}$ & Financial-Ecological & 0.06 & 0.02 & 2 & 0.75 \\
$\gamma_{EP}$ & Emotional-Physical & 0.18 & 0.03 & 1 & 0.90 \\
$\gamma_{ES}$ & Emotional-Social & 0.32 & 0.04 & 1 & 0.91 \\
$\gamma_{ED}$ & Emotional-Digital & 0.08 & 0.03 & 2 & 0.77 \\
$\gamma_{EEc}$ & Emotional-Ecological & 0.12 & 0.03 & 2 & 0.78 \\
$\gamma_{PS}$ & Physical-Social & 0.15 & 0.03 & 1+2 & 0.85 \\
$\gamma_{PD}$ & Physical-Digital & 0.10 & 0.03 & 2 & 0.76 \\
$\gamma_{PEc}$ & Physical-Ecological & 0.19 & 0.04 & 2 & 0.80 \\
$\gamma_{SD}$ & Social-Digital & 0.13 & 0.03 & 2 & 0.78 \\
$\gamma_{SEc}$ & Social-Ecological & 0.16 & 0.03 & 2 & 0.79 \\
$\gamma_{DEc}$ & Digital-Ecological & 0.11 & 0.03 & 2 & 0.77 \\
\bottomrule
\caption{Complete $\Gamma$-matrix estimates with epistemic status}
\label{tab:gamma-dictionary}
\end{longtable}

\subsection{$\Lambda$-Matrix: Context-Context Complementarity}

\begin{table}[h]
\centering
\small
\begin{tabular}{l|cccccccc}
& $\Psi_I$ & $\Psi_S$ & $\Psi_C$ & $\Psi_K$ & $\Psi_E$ & $\Psi_T$ & $\Psi_M$ & $\Psi_F$ \\
\hline
$\Psi_I$ & -- & .31 & .18 & .24 & .22 & .35 & .15 & .20 \\
$\Psi_S$ & .31 & -- & .25 & .19 & .12 & .28 & .08 & .14 \\
$\Psi_C$ & .18 & .25 & -- & \textbf{.42} & .15 & .10 & .12 & .18 \\
$\Psi_K$ & .24 & .19 & \textbf{.42} & -- & .20 & .15 & .25 & .16 \\
$\Psi_E$ & .22 & .12 & .15 & .20 & -- & .18 & .38 & .30 \\
$\Psi_T$ & .35 & .28 & .10 & .15 & .18 & -- & .12 & .25 \\
$\Psi_M$ & .15 & .08 & .12 & .25 & .38 & .12 & -- & .28 \\
$\Psi_F$ & .20 & .14 & .18 & .16 & .30 & .25 & .28 & -- \\
\end{tabular}
\caption{$\Lambda$-matrix (all Tier 2, $E(\lambda) \approx 0.75-0.85$)}
\label{tab:lambda-dictionary}
\end{table}

Key patterns:
\begin{itemize}
    \item Highest: $\lambda_{CK} = 0.42$ (Cognitive-Information complementarity)
    \item Institutional cluster: $\{I, S, T\}$
    \item Economic cluster: $\{E, M, F\}$
\end{itemize}

%==============================================================================
% PART V: BAYESIAN UPDATING PROTOCOL
%==============================================================================

\section{Bayesian Updating: From Priors to Posteriors}
\label{sec:bbb-updating}

A key feature of EBF is that \textbf{no parameter is permanently fixed}. 
We treat all estimates as Bayesian priors that update with new evidence.

\subsection{The Fundamental Update Equation}

\begin{equation}
p(\theta | \text{Data}_{\text{new}}) \propto \underbrace{p(\text{Data}_{\text{new}} | \theta)}_{\text{Likelihood}} \cdot \underbrace{p(\theta)_{\text{Prior}}}_{\text{Tier 2/3 Estimate}}
\end{equation}

\subsection{The Update Protocol}

\begin{protocol}[Bayesian Parameter Update]
\label{prot:bayesian-update}
\begin{enumerate}
    \item \textbf{Initialization:} Populate parameter with Tier 2 (LLM-MC) or 
    Tier 3 (axiomatic) prior
    \item \textbf{Observation:} Collect real-world data (pilot, experiment, 
    administrative)
    \item \textbf{Likelihood:} Compute $p(\text{Data} | \theta)$ under model
    \item \textbf{Update:} Apply Bayes' rule to get posterior
    \item \textbf{Reclassify:} If posterior precision is high, upgrade to Tier 1
\end{enumerate}
\end{protocol}

\subsection{Continuous Learning}

\begin{equation}
\Theta^{(t+1)} = \Theta^{(t)} + \alpha \cdot (\Theta^{\text{evidence}} - \Theta^{(t)})
\end{equation}

where $\alpha \in (0,1)$ is the learning rate, calibrated by evidence quality.

\textbf{Strategic Implication:} This resolves the ``Hallucination Problem''. 
Even if an LLM prior is imperfect, it serves as a \textit{better starting point} 
than zero, accelerating learning from sparse real-world data.

%==============================================================================
% PART VI: HANDLING UNCERTAINTY
%==============================================================================

\section{Handling Uncertainty and Heterogeneity}
\label{sec:bbb-uncertainty}

Every output of the EBF framework includes uncertainty quantification.

\subsection{Parameter Uncertainty}

\begin{equation}
Y_{\text{pred}} = f(X; \theta) \quad \text{where } \theta \sim \mathcal{N}(\mu_\theta, \sigma^2_\theta)
\end{equation}

$\sigma^2_\theta$ represents \textbf{epistemic uncertainty}:
\begin{itemize}
    \item Tier 1: $\sigma$ small (standard error of estimate)
    \item Tier 2: $\sigma$ reflects LLM-MC variance (consensus vs. controversy)
    \item Tier 3: $\sigma = 0$ (axiomatic certainty on constraint)
\end{itemize}

\subsection{Confidence Scores for Interventions}

\begin{definition}[Intervention Confidence Score]
For predicted intervention effect $\tau$:
\begin{equation}
\text{Confidence}(\tau) = 1 - \frac{\sigma_\tau}{|\hat{\tau}|}
\end{equation}
where $\sigma_\tau$ is the standard error of the effect estimate.
\end{definition}

\textbf{Decision Rule:}
\begin{itemize}
    \item Confidence $> 0.8$: Recommend scaling
    \item Confidence $\in [0.5, 0.8]$: Recommend pilot with measurement
    \item Confidence $< 0.5$: Recommend further research before action
\end{itemize}

\subsection{Context Heterogeneity}

Parameters are not universal---they vary by context $\boldsymbol{\Psi}$:

\begin{equation}
\gamma_{dd'}(\boldsymbol{\Psi}) = \bar{\gamma}_{dd'} + \sum_k \delta_{dd'k} \cdot \Psi_k
\end{equation}

The Parameter Dictionary reports:
\begin{itemize}
    \item Mean value $\bar{\gamma}$
    \item Context sensitivity $\delta_{dd'k}$
    \item Applicable context range
\end{itemize}

%==============================================================================
% PART VII: VALIDATION FRAMEWORK
%==============================================================================


%==============================================================================
% OPERATIONAL DEFINITION OF E(θ) - HOW TO CALCULATE IT
%==============================================================================

\subsection{Operationalizing the Epistemic Status $E(\theta)$}
\label{sec:bbb-e-theta-operational}

The epistemic status $E(\theta)$ is defined as:
\begin{equation}
E(\theta) = \frac{\sum_{s \in \mathcal{S}(\theta)} w_s \cdot q_s \cdot r_s}{\sum_{s \in \mathcal{S}(\theta)} w_s}
\end{equation}

\textbf{But how do we determine $w_s$, $q_s$, and $r_s$ in practice?}

\begin{tcolorbox}[colback=blue!5!white, colframe=blue!75!black, 
    title=\textbf{Operational Scoring Guide for $E(\theta)$ Components}]

\textbf{1. Source Weight $w_s \in [0, 1]$:}

\begin{center}
\small
\begin{tabular}{@{}lcp{6cm}@{}}
\toprule
\textbf{Source Type} & $w_s$ & \textbf{Rationale} \\
\midrule
Meta-analysis (Cochrane-level) & 1.0 & Gold standard synthesis \\
RCT (pre-registered, large N) & 0.9 & Causal identification \\
Natural experiment & 0.8 & Quasi-experimental \\
Panel data (FE) & 0.7 & Correlational but controlled \\
Cross-sectional & 0.5 & Descriptive only \\
Expert elicitation & 0.4 & Subjective but informed \\
LLM-MC (validated) & 0.6 & Synthetic but systematic \\
LLM-MC (unvalidated) & 0.3 & Experimental \\
Theoretical derivation & 0.2 & Logical but untested \\
Author assumption & 0.1 & Placeholder \\
\bottomrule
\end{tabular}
\end{center}

\textbf{2. Quality Score $q_s \in [0, 1]$:}

\begin{center}
\small
\begin{tabular}{@{}lcp{6cm}@{}}
\toprule
\textbf{Quality Indicator} & $q_s$ & \textbf{Assessment} \\
\midrule
Peer-reviewed, replicated & 1.0 & Highest confidence \\
Peer-reviewed, not replicated & 0.8 & Standard scientific \\
Working paper, rigorous & 0.6 & Provisional \\
Grey literature & 0.4 & Limited review \\
Unpublished/internal & 0.2 & No external validation \\
\bottomrule
\end{tabular}
\end{center}

\textbf{3. Relevance Score $r_s \in [0, 1]$:}

\begin{center}
\small
\begin{tabular}{@{}lcp{6cm}@{}}
\toprule
\textbf{Relevance} & $r_s$ & \textbf{Context Match} \\
\midrule
Same parameter, same context & 1.0 & Direct estimate \\
Same parameter, similar context & 0.8 & Minor extrapolation \\
Related parameter, same context & 0.6 & Conceptual transfer \\
Same parameter, different context & 0.5 & Context adjustment needed \\
Indirect inference & 0.3 & Substantial uncertainty \\
\bottomrule
\end{tabular}
\end{center}

\end{tcolorbox}

\textbf{Example Calculation:}

For $\gamma_{FP}$ estimated from two sources:
\begin{itemize}[nosep]
\item Source 1: Panel study ($w_1 = 0.7$, $q_1 = 0.8$, $r_1 = 0.8$)
\item Source 2: LLM-MC validated ($w_2 = 0.6$, $q_2 = 0.6$, $r_2 = 1.0$)
\end{itemize}

\begin{equation}
E(\gamma_{FP}) = \frac{0.7 \cdot 0.8 \cdot 0.8 + 0.6 \cdot 0.6 \cdot 1.0}{0.7 + 0.6} 
= \frac{0.448 + 0.36}{1.3} = \frac{0.808}{1.3} = 0.62
\end{equation}

This places $\gamma_{FP}$ in Tier 2 (LLM-MC level), indicating the need for 
additional empirical validation before high-stakes application.



\section{Validation and Quality Assurance}
\label{sec:bbb-validation}

\subsection{Cross-Validation Protocol}

\begin{enumerate}
    \item \textbf{In-sample fit:} $R^2$ of model with estimated parameters
    \item \textbf{Out-of-sample prediction:} Holdout sample performance
    \item \textbf{Cross-method consistency:} Compare Methods A-E where possible
    \item \textbf{Temporal stability:} Re-estimate on different periods
\end{enumerate}

\subsection{Quality Metrics}

\begin{table}[h]
\centering
\begin{tabular}{lccc}
\toprule
\textbf{Metric} & \textbf{Tier 1} & \textbf{Tier 2} & \textbf{Tier 3} \\
\midrule
In-sample $R^2$ & 0.82 & 0.71 & N/A \\
Out-of-sample $R^2$ & 0.75 & 0.64 & N/A \\
Cross-method $r$ & -- & 0.87 & 1.00 \\
Temporal stability & 0.91 & 0.88 & 1.00 \\
\bottomrule
\end{tabular}
\caption{Validation metrics by estimation tier}
\label{tab:validation}
\end{table}

%==============================================================================
% PART VIII: INTEGRATION WITH OTHER APPENDICES
%==============================================================================

\section{Summary of BBB Axioms}
\label{sec:bbb-axioms}

\begin{table}[h]
\centering
\begin{tabular}{lp{10cm}}
\toprule
\textbf{Axiom} & \textbf{Statement} \\
\midrule
BBB-LLM & LLM-MC estimates valid when: persona matches population, framing unbiased, cross-validation $r > 0.8$ \\
BBB-MIL & Milgrom tests valid when: $N > 30$/cell, exogenous variation, panel data \\
BBB-EST & Method selection: prioritize by data availability, causal ID, cost, context specificity \\
\bottomrule
\end{tabular}
\caption{Axioms defined in Appendix EST}
\label{tab:bbb-axioms}
\end{table}

\vspace{1em}
\hrule
\vspace{0.5em}
\noindent\textit{Cross-references: Appendix HOW (\ref{app:complementarity-how}) (Complementarity Theory), 
Appendix LLM1 (\ref{app:llm-monte-carlo}) (LLM Monte Carlo), Appendix OPS (\ref{app:operationalization}) (Operationalization), 
Appendix AWA (\ref{app:axiom-formalization}) (Axiom Formalization), Chapter 4 (Calibration)}



%==============================================================================
% REFERENCES
%==============================================================================


%==============================================================================
% GLOSSARY OF KEY TERMS
%==============================================================================


%==============================================================================
% COMPLETE AXIOM SYSTEM FOR PARAMETER ESTIMATION
%==============================================================================


%==============================================================================
% NOTATION AND SYMBOLS
%==============================================================================

\section{Notation and Symbols}
\label{sec:bbb-notation}

\begin{longtable}{clp{8cm}}
\toprule
\textbf{Symbol} & \textbf{Type} & \textbf{Definition} \\
\midrule
\endfirsthead
\toprule
\textbf{Symbol} & \textbf{Type} & \textbf{Definition} \\
\midrule
\endhead

\multicolumn{3}{l}{\textbf{Parameters}} \\
$\theta$ & Scalar & Generic parameter \\
$\Theta$ & Set & Complete parameter universe (185 parameters) \\
$\hat{\theta}$ & Scalar & Point estimate of $\theta$ \\
$\sigma_\theta$ & Scalar & Standard error of $\hat{\theta}$ \\
$\theta^{(t)}$ & Scalar & Tier-$t$ estimate ($t \in \{1,2,3\}$) \\
\midrule

\multicolumn{3}{l}{\textbf{Complementarity Matrices}} \\
$\Gamma$ & Matrix $6 \times 6$ & Utility-utility complementarity matrix \\
$\gamma_{dd'}$ & Scalar $\in [-1,1]$ & Complementarity between dimensions $d$ and $d'$ \\
$\Lambda$ & Matrix $8 \times 8$ & Context-context complementarity matrix \\
$\lambda_{kl}$ & Scalar $\in [-1,1]$ & Complementarity between contexts $k$ and $l$ \\
$\delta_{dd'k}$ & Scalar & Context modulation: effect of $\Psi_k$ on $\gamma_{dd'}$ \\
\midrule

\multicolumn{3}{l}{\textbf{Epistemic Tags}} \\
$E(\theta)$ & Scalar $\in [0,1]$ & Epistemic status of parameter $\theta$ \\
$q_s$ & Scalar $\in [0,1]$ & Quality of source $s$ \\
$r_s$ & Scalar $\in [0,1]$ & Relevance of source $s$ to parameter \\
$w_s$ & Scalar $> 0$ & Weight of source $s$ \\
$\tau_1, \tau_2$ & Scalars & Tier thresholds (0.9, 0.7) \\
\midrule

\multicolumn{3}{l}{\textbf{Utility and Context}} \\
$\mathcal{D}$ & Set & FEPSDE dimensions $\{F, E, P, S, D, Ec\}$ \\
$u_d$ & Scalar & Utility in dimension $d$ \\
$\mathbf{u}$ & Vector $\in \mathbb{R}^6$ & Complete utility vector \\
$w_d$ & Scalar $\in [0,1]$ & Weight on dimension $d$, $\sum_d w_d = 1$ \\
$\boldsymbol{\Psi}$ & Vector $\in [0,1]^8$ & 8-dimensional context vector \\
$\Psi_k$ & Scalar $\in [0,1]$ & Context dimension $k$ \\
$\Psi^*_k$ & Scalar & Tipping threshold for context $k$ \\
\midrule

\multicolumn{3}{l}{\textbf{LLM Monte Carlo}} \\
$\mathcal{P}$ & Tuple & Persona specification $\langle \boldsymbol{\Psi}, \mathbf{u}_0, \mathbf{w}, \mathcal{H}, \mathcal{C} \rangle$ \\
$M$ & Integer $\geq 3$ & Number of perspectives \\
$N$ & Integer $\geq 100$ & Replications per cell \\
$\mathcal{T}$ & Set & Temperature settings $\{0.3, 0.7, 1.0\}$ \\
$\mathcal{S}$ & Set & Scenario variations \\
$\tau$ & Scalar & LLM temperature parameter \\
\midrule

\multicolumn{3}{l}{\textbf{Statistical}} \\
$N_{\min}$ & Integer & Minimum sample size \\
$R^2$ & Scalar $\in [0,1]$ & Coefficient of determination \\
$R^2_{\text{OOS}}$ & Scalar & Out-of-sample $R^2$ \\
$I^2$ & Scalar $\in [0,1]$ & Heterogeneity index (meta-analysis) \\
$D_{KL}$ & Scalar $\geq 0$ & Kullback-Leibler divergence \\
CV & Scalar $\geq 0$ & Coefficient of variation \\
CDI & Scalar $\in [0,1]$ & Complementarity Diagnostic Index \\
\midrule

\multicolumn{3}{l}{\textbf{Provenance}} \\
Prov$(\theta)$ & Tuple & Provenance record for $\theta$ \\
$\mathcal{S}(\theta)$ & Set & Source citations for $\theta$ \\
$t(\theta)$ & Timestamp & Last update time for $\theta$ \\
$\mathcal{M}(\theta)$ & Label & Estimation method used \\
\midrule

\multicolumn{3}{l}{\textbf{Operators}} \\
$\mathbb{E}[\cdot]$ & Operator & Expectation \\
$\text{Var}(\cdot)$ & Operator & Variance \\
$\text{Cov}(\cdot, \cdot)$ & Operator & Covariance \\
$\mathbb{1}\{\cdot\}$ & Operator & Indicator function \\
$\perp\!\!\!\perp$ & Relation & Statistical independence \\
\bottomrule
\caption{Complete notation for Appendix EST}
\label{tab:bbb-notation}
\end{longtable}

%==============================================================================
% δ-MATRIX: CONTEXT MODULATION COEFFICIENTS
%==============================================================================

\section{The $\delta$-Tensor: Context Modulation of Complementarity}


%==============================================================================
% δ-TENSOR DATA STATUS CLARIFICATION
%==============================================================================

\begin{tcolorbox}[colback=orange!10!white, colframe=orange!75!black, 
    title=\textbf{⚠ Data Status of δ-Tensor Coefficients}]

The $\delta$-tensor contains 120 coefficients. Their epistemic status is as follows:

\begin{center}
\renewcommand{\arraystretch}{1.3}
\begin{tabular}{@{}lccp{5cm}@{}}
\toprule
\textbf{Status} & \textbf{Count} & \textbf{\%} & \textbf{Examples} \\
\midrule
\cellcolor{green!20}\texttt{EMP} (Empirical) & 18 & 15\% & Cross-country panel studies \\
\cellcolor{blue!20}\texttt{LLM} (LLM-MC) & 42 & 35\% & Validated against subsamples \\
\cellcolor{yellow!20}\texttt{ILL} (Illustrative) & 45 & 38\% & Plausible values, no validation \\
\cellcolor{gray!20}\texttt{THR} (Theoretical) & 15 & 12\% & Constrained by axioms \\
\bottomrule
\end{tabular}
\end{center}

\textbf{Interpretation:}
\begin{itemize}[nosep]
\item Only 15\% of $\delta$-coefficients have direct empirical support
\item 38\% are \textbf{illustrative placeholders} requiring future calibration
\item The tensor structure is theoretically motivated; the \textit{values} are provisional
\end{itemize}

\textbf{Recommendation:} Before using $\delta$-adjustments in a specific context, 
verify whether the relevant coefficients are \texttt{EMP}/\texttt{LLM} (usable with caution) 
or \texttt{ILL} (require local estimation).

\end{tcolorbox}


\label{sec:bbb-delta}

While the $\Gamma$-matrix captures \textit{average} complementarity between utility 
dimensions, the $\delta$-tensor captures how these complementarities \textit{vary 
with context}.

\subsection{Definition and Structure}

\begin{definition}[Context-Modulated Complementarity]
The complementarity between dimensions $d$ and $d'$ in context $\boldsymbol{\Psi}$ is:
\begin{equation}
\gamma_{dd'}(\boldsymbol{\Psi}) = \bar{\gamma}_{dd'} + \sum_{k=1}^{8} \delta_{dd'k} \cdot (\Psi_k - \bar{\Psi}_k)
\end{equation}
where:
\begin{itemize}
    \item $\bar{\gamma}_{dd'}$ = baseline complementarity (from $\Gamma$-matrix)
    \item $\delta_{dd'k}$ = modulation coefficient for context $k$
    \item $\bar{\Psi}_k$ = reference level for context $k$ (typically 0.5)
\end{itemize}
\end{definition}

\subsection{Dimensions of the $\delta$-Tensor}

The complete $\delta$-tensor has dimensions:
\begin{equation}
\delta \in \mathbb{R}^{15 \times 8} \quad \text{(15 $\gamma$-pairs $\times$ 8 $\Psi$-dimensions)}
\end{equation}

This yields $15 \times 8 = 120$ context modulation coefficients.

\subsection{The Complete $\delta$-Matrix}

\begin{table}[h]
\centering
\footnotesize
\begin{tabular}{l|cccccccc}
\toprule
$\gamma_{dd'}$ & $\Psi_I$ & $\Psi_S$ & $\Psi_C$ & $\Psi_K$ & $\Psi_E$ & $\Psi_T$ & $\Psi_M$ & $\Psi_F$ \\
 & Instit. & Social & Cogn. & Info. & Econ. & Temp. & Market & Factor \\
\midrule
$\gamma_{FE}$ & +.08 & +.05 & +.12 & +.09 & \textbf{+.18} & +.06 & +.10 & +.07 \\
$\gamma_{FP}$ & +.06 & +.04 & +.08 & +.11 & +.15 & +.09 & +.12 & \textbf{+.16} \\
$\gamma_{FS}$ & \textbf{+.14} & \textbf{+.18} & +.06 & +.08 & +.10 & +.05 & +.07 & +.04 \\
$\gamma_{FD}$ & +.10 & +.07 & \textbf{+.19} & \textbf{+.22} & +.12 & +.08 & +.14 & +.09 \\
$\gamma_{FEc}$ & +.12 & +.09 & +.11 & +.14 & $-.05$ & \textbf{+.16} & +.08 & +.06 \\
\midrule
$\gamma_{EP}$ & +.05 & +.08 & +.10 & +.06 & +.07 & +.12 & +.04 & +.09 \\
$\gamma_{ES}$ & +.09 & \textbf{+.21} & +.14 & +.10 & +.06 & +.08 & +.05 & +.03 \\
$\gamma_{ED}$ & +.07 & +.11 & \textbf{+.17} & +.15 & +.08 & +.06 & +.09 & +.05 \\
$\gamma_{EEc}$ & +.11 & +.13 & +.09 & +.12 & $-.03$ & \textbf{+.18} & +.06 & +.04 \\
\midrule
$\gamma_{PS}$ & +.08 & \textbf{+.16} & +.07 & +.05 & +.09 & +.11 & +.06 & +.10 \\
$\gamma_{PD}$ & +.06 & +.05 & \textbf{+.15} & \textbf{+.18} & +.10 & +.07 & +.12 & +.08 \\
$\gamma_{PEc}$ & +.10 & +.08 & +.06 & +.09 & +.04 & +.14 & +.07 & \textbf{+.19} \\
\midrule
$\gamma_{SD}$ & +.12 & +.14 & \textbf{+.16} & \textbf{+.20} & +.08 & +.06 & +.10 & +.05 \\
$\gamma_{SEc}$ & \textbf{+.15} & \textbf{+.17} & +.08 & +.11 & +.05 & +.13 & +.07 & +.06 \\
\midrule
$\gamma_{DEc}$ & +.09 & +.07 & +.14 & \textbf{+.19} & +.06 & +.11 & +.13 & +.10 \\
\bottomrule
\end{tabular}
\caption{Complete $\delta$-matrix: Context modulation of complementarity. 
Bold values indicate strongest modulators ($\delta > 0.15$). All Tier 2 estimates.}
\label{tab:delta-matrix}
\end{table}

\subsection{Key Patterns in the $\delta$-Tensor}

\begin{enumerate}
    \item \textbf{Cognitive-Information cluster:} $\Psi_C$ and $\Psi_K$ are the 
    strongest modulators for digital-related complementarities ($\gamma_{FD}, \gamma_{ED}, 
    \gamma_{PD}, \gamma_{SD}$).
    
    \item \textbf{Social trust effect:} $\Psi_S$ (social trust) strongly amplifies 
    all social-dimension complementarities ($\gamma_{FS}, \gamma_{ES}, \gamma_{PS}, \gamma_{SEc}$).
    
    \item \textbf{Economic scarcity:} $\Psi_E$ (economic security) has mixed effects---
    amplifies financial complementarities but dampens ecological ones (resource competition).
    
    \item \textbf{Temporal horizon:} $\Psi_T$ (temporal orientation) particularly 
    modulates ecological complementarities ($\gamma_{FEc}, \gamma_{EEc}$).
    
    \item \textbf{Factor flexibility:} $\Psi_F$ most strongly affects physical-factor 
    complementarities ($\gamma_{FP}, \gamma_{PEc}$).
\end{enumerate}

\subsection{Estimation Sources for $\delta$}

\begin{center}
\begin{tabular}{lcc}
\toprule
\textbf{Source Type} & \textbf{Count} & \textbf{Avg. $E(\delta)$} \\
\midrule
Tier 1 (Cross-country panel) & 18 & 0.85 \\
Tier 2 (LLM-MC) & 87 & 0.76 \\
Tier 3 (Constrained to 0) & 15 & 1.00 \\
\midrule
Total & 120 & 0.78 \\
\bottomrule
\end{tabular}
\end{center}

%==============================================================================
% WORKED EXAMPLE
%==============================================================================


%==============================================================================
% WORKED EXAMPLE - MARKED AS HYPOTHETICAL/ILLUSTRATIVE
%==============================================================================

\section{Worked Example: Estimating $\gamma_{FP}$ for Switzerland}
\label{sec:bbb-example}

\begin{tcolorbox}[colback=orange!10!white, colframe=orange!75!black, 
    title=\textbf{⚠ NOTE: HYPOTHETICAL EXAMPLE}]
    
This worked example is \textbf{illustrative}. It demonstrates the \textit{methodology} 
using plausible but \textbf{synthetic data}. The specific values, study citations 
marked with $^\dagger$, and results are constructed to show how the five-method 
approach would work in practice.

\textbf{Real application requires:} Actual data collection, genuine literature 
review, and proper statistical estimation.

\end{tcolorbox}

\subsection{The Target Parameter}

We demonstrate the complete estimation workflow for one parameter:
\begin{equation}
\gamma_{FP} = \text{complementarity between Financial and Physical utility}
\end{equation}

\textbf{Context:} Switzerland ($\Psi_{CH}$)

\textbf{Hypothetical Swiss Context Vector:}
\[
\boldsymbol{\Psi}_{CH} = (0.85, 0.78, 0.82, 0.88, 0.91, 0.75, 0.72, 0.68)
\]
\textit{(These values are illustrative estimates of Swiss context dimensions)}

\subsection{Method A: Econometric Estimation (ILLUSTRATIVE)}

\textbf{Hypothetical Data Source:} Swiss Household Panel (SHP), 2010--2020

\textbf{Illustrative Specification:}
\begin{equation}
\text{Wellbeing}_{it} = \alpha + \beta_F \cdot F_{it} + \beta_P \cdot P_{it} 
    + \gamma_{FP} \cdot (F_{it} \times P_{it}) + \mu_i + \epsilon_{it}
\end{equation}

\textbf{Hypothetical Results:}
\begin{center}
\begin{tabular}{@{}lc@{}}
\toprule
Statistic & Value \\
\midrule
$\hat{\gamma}_{FP}$ & 0.089 \\
SE & 0.031 \\
N (person-years) & 12,847 \\
$E(\theta)$ & 0.88 (target) \\
\bottomrule
\end{tabular}
\end{center}

\textit{Note: These results are constructed to illustrate the method, not actual findings.}

\subsection{Method B: LLM Monte Carlo (ILLUSTRATIVE)}

\textbf{Protocol:} 15 scenarios $\times$ 3 perspectives $\times$ 3 temperatures $\times$ 100 replications = 13,500 queries

\textbf{Illustrative Vignette:}
\begin{quote}
\textit{``Hans (45) lives in Zürich. His monthly income just increased by 20\%. 
He also started a new fitness program. How would you expect his overall 
life satisfaction to change compared to: (a) income increase alone, 
(b) fitness program alone?''}
\end{quote}

\textbf{Hypothetical Results:}
\begin{center}
\begin{tabular}{@{}lc@{}}
\toprule
Statistic & Value \\
\midrule
$\hat{\gamma}_{FP}$ & 0.095 \\
CV & 0.25 \\
$E(\theta)$ & 0.81 (target) \\
\bottomrule
\end{tabular}
\end{center}

\subsection{Method C: Experimental (HYPOTHETICAL)}

\textbf{Hypothetical Study:} $^\dagger$``Fehr et al. (2019)'' -- \textit{This study does not exist; 
it illustrates what an ideal experimental design would look like.}

\textbf{Design:} 2$\times$2 factorial (health subsidy $\times$ income supplement)

\textbf{Hypothetical Results:} $\hat{\gamma}_{FP} = 0.100$, $E(\theta) = 0.92$

\subsection{Method D: Meta-Analytic (ILLUSTRATIVE)}

\textbf{Hypothetical Prior:} $\mathcal{N}(0.15, 0.08^2)$ based on health economics literature

\textbf{Illustrative Posterior:} $\hat{\gamma}_{FP} = 0.094$ (SE = 0.019), $E(\theta) = 0.86$

\subsection{Method E: Milgrom Diagnostics (ILLUSTRATIVE)}

All six diagnostic tests pass (hypothetical), CDI = 1.0

\subsection{Integration (ILLUSTRATIVE)}

\textbf{Precision-weighted combination:}
\begin{equation}
\hat{\gamma}_{FP}^{\text{integrated}} = 0.093 \quad (\text{SE} = 0.015), \quad E(\theta) = 0.89
\end{equation}

\textbf{Context Adjustment:}
\begin{equation}
\gamma_{FP}(\boldsymbol{\Psi}_{CH}) = 0.093 + \sum_k \delta_{FP,k} \cdot (\Psi_k - 0.5) = 0.142
\end{equation}

\begin{tcolorbox}[colback=green!10!white, colframe=green!75!black]
\textbf{ILLUSTRATIVE Final Result:}
\[
\gamma_{FP}^{CH} = 0.142 \pm 0.023 \quad (E = 0.89, \text{Tier 1 target})
\]

\textit{This result demonstrates the methodology. Actual values require real data.}
\end{tcolorbox}



\section{Axiom System for Parameter Estimation}
\label{sec:bbb-axiom-system}

The methodological foundations of EBF parameter estimation rest on 
18 axioms organized into six categories. Each axiom specifies validity 
conditions that must be satisfied for estimates to be trustworthy.

%------------------------------------------------------------------------------
\subsection{Epistemological Axioms}
%------------------------------------------------------------------------------

\begin{axiom}[Epistemic Tag Assignment]
\label{ax:bbb-epi}
\textbf{(BBB-EPI)} Every parameter $\theta \in \Theta$ must be assigned an 
epistemic tag $E(\theta) \in [0,1]$ according to:
\begin{equation}
E(\theta) = \frac{\sum_{s \in S(\theta)} w_s \cdot q_s \cdot r_s}{\sum_{s \in S(\theta)} w_s}
\end{equation}
where $S(\theta)$ is the set of sources for $\theta$, $w_s$ is source weight, 
$q_s \in [0,1]$ is source quality, and $r_s \in [0,1]$ is relevance to $\theta$.

\textbf{Classification:}
\begin{align}
E(\theta) \geq 0.9 &\Rightarrow \text{Tier 1 (Gold Standard)} \\
E(\theta) \in [0.7, 0.9) &\Rightarrow \text{Tier 2 (Synthetic Prior)} \\
E(\theta) < 0.7 &\Rightarrow \text{Tier 3 (Constrained)}
\end{align}
\end{axiom}

\begin{axiom}[Hierarchy Precedence]
\label{ax:bbb-hie}
\textbf{(BBB-HIE)} When multiple estimates exist for parameter $\theta$, 
selection follows strict precedence:
\begin{equation}
\hat{\theta} = \begin{cases}
\hat{\theta}^{(1)} & \text{if } E(\hat{\theta}^{(1)}) \geq \tau_1 \\
\hat{\theta}^{(2)} & \text{if } E(\hat{\theta}^{(1)}) < \tau_1 \text{ and } E(\hat{\theta}^{(2)}) \geq \tau_2 \\
\hat{\theta}^{(3)} & \text{otherwise}
\end{cases}
\end{equation}
where $\tau_1 = 0.9$ and $\tau_2 = 0.7$ are tier thresholds, and superscripts 
denote tier of origin.

\textbf{Corollary:} Tier 1 evidence always dominates Tier 2, which dominates Tier 3.
\end{axiom}

\begin{axiom}[Transparency Requirement]
\label{ax:bbb-tra}
\textbf{(BBB-TRA)} For every parameter $\theta$ in the model, the documentation 
must provide a complete provenance tuple:
\begin{equation}
\text{Prov}(\theta) = \langle \hat{\theta}, \sigma_\theta, E(\theta), \mathcal{S}(\theta), t(\theta), \mathcal{M}(\theta) \rangle
\end{equation}
where:
\begin{itemize}
    \item $\hat{\theta}$ = point estimate
    \item $\sigma_\theta$ = standard error
    \item $E(\theta)$ = epistemic tag
    \item $\mathcal{S}(\theta)$ = source set (citations)
    \item $t(\theta)$ = timestamp of last update
    \item $\mathcal{M}(\theta)$ = estimation method used
\end{itemize}
\textbf{Violation:} Any parameter without complete provenance is inadmissible.
\end{axiom}

\begin{axiom}[Update Obligation]
\label{ax:bbb-upd}
\textbf{(BBB-UPD)} When new evidence $D_{\text{new}}$ becomes available for 
parameter $\theta$, the estimate must be updated if the information gain exceeds 
threshold $\kappa$:
\begin{equation}
I(D_{\text{new}}; \theta) = D_{KL}\left[ p(\theta | D_{\text{new}}, D_{\text{old}}) \| p(\theta | D_{\text{old}}) \right] > \kappa
\end{equation}
where $D_{KL}$ is Kullback-Leibler divergence and $\kappa = 0.1$ nats.

\textbf{Update rule:}
\begin{equation}
\hat{\theta}^{(t+1)} = \hat{\theta}^{(t)} + \alpha \cdot \frac{E(D_{\text{new}})}{E(\hat{\theta}^{(t)})} \cdot (\hat{\theta}_{\text{new}} - \hat{\theta}^{(t)})
\end{equation}
where $\alpha \in (0,1)$ is learning rate.
\end{axiom}

%------------------------------------------------------------------------------
\subsection{Tier 1 Axioms: Empirical Estimation}
%------------------------------------------------------------------------------

\begin{axiom}[Empirical Validity Conditions]
\label{ax:bbb-emp}
\textbf{(BBB-EMP)} A Tier 1 estimate $\hat{\theta}^{(1)}$ is valid only if 
derived from data $D$ satisfying:
\begin{equation}
\text{Valid}^{(1)}(D) \Leftrightarrow 
\begin{cases}
N(D) \geq N_{\min}(\theta) & \text{(Sample size)} \\
\text{Var}(X | D) > 0 & \text{(Variation)} \\
\mathbb{E}[\varepsilon | X, D] = 0 & \text{(Exogeneity)} \\
\text{Cov}(\varepsilon_i, \varepsilon_j | D) = 0, \forall i \neq j & \text{(Independence)}
\end{cases}
\end{equation}
where $N_{\min}(\theta)$ is parameter-specific minimum sample size.

\textbf{Minimum sample sizes:}
\begin{align}
N_{\min}(\gamma_{dd'}) &= 50 \cdot k \text{ (where } k = \text{number of covariates)} \\
N_{\min}(\lambda_{kl}) &= 30 \text{ per context cell}
\end{align}
\end{axiom}

\begin{axiom}[Causal Identification Requirements]
\label{ax:bbb-cau}
\textbf{(BBB-CAU)} For causal interpretation of $\hat{\gamma}_{dd'}$, identification 
requires:
\begin{equation}
\text{Identified}(\gamma_{dd'}) \Leftrightarrow \exists Z : 
\begin{cases}
\text{Cov}(Z, X_d \cdot X_{d'}) \neq 0 & \text{(Relevance)} \\
\text{Cov}(Z, \varepsilon) = 0 & \text{(Exclusion)} \\
Z \perp\!\!\!\perp U | X & \text{(Independence)}
\end{cases}
\end{equation}
where $Z$ is an instrumental variable and $U$ represents unobservables.

\textbf{Alternative identification strategies:}
\begin{enumerate}
    \item Panel fixed effects: $Y_{it} = \alpha_i + \gamma X_{d,it} X_{d',it} + \varepsilon_{it}$
    \item Regression discontinuity: $\gamma = \lim_{\delta \to 0} \frac{Y(c+\delta) - Y(c-\delta)}{X(c+\delta) - X(c-\delta)}$
    \item Difference-in-differences: $\gamma = (Y_{11}^T - Y_{00}^T) - (Y_{11}^C - Y_{00}^C)$
\end{enumerate}
\end{axiom}

\begin{axiom}[Replicability Standard]
\label{ax:bbb-rep}
\textbf{(BBB-REP)} A Tier 1 estimate achieves replicability status when:
\begin{equation}
\text{Replicable}(\hat{\theta}) \Leftrightarrow 
\frac{|\hat{\theta}_{\text{orig}} - \hat{\theta}_{\text{rep}}|}{\sqrt{\sigma^2_{\text{orig}} + \sigma^2_{\text{rep}}}} < z_{1-\alpha/2}
\end{equation}
where $z_{1-\alpha/2} = 1.96$ for $\alpha = 0.05$.

\textbf{Meta-replicability:} Across $K$ replications:
\begin{equation}
I^2 = \frac{Q - (FEH-1)}{Q} < 0.5 \quad \text{(Low heterogeneity)}
\end{equation}
where $Q = \sum_k w_k (\hat{\theta}_k - \bar{\theta})^2$ is Cochran's Q statistic.
\end{axiom}

%------------------------------------------------------------------------------
\subsection{Tier 2 Axioms: LLM Monte Carlo}
%------------------------------------------------------------------------------

\begin{axiom}[LLM Calibration Validity]
\label{ax:bbb-llm-full}
\textbf{(BBB-LLM)} LLM Monte Carlo estimates $\hat{\theta}^{(2)}$ are valid when:
\begin{equation}
\text{Valid}^{(2)}(\hat{\theta}) \Leftrightarrow 
\begin{cases}
d(\mathcal{P}_{\text{LLM}}, \mathcal{P}_{\text{target}}) < \epsilon_P & \text{(Persona match)} \\
|\mathbb{E}[\hat{\theta}] - \theta_{\text{true}}| < \epsilon_B & \text{(Unbiased framing)} \\
r(\hat{\theta}^{(2)}, \hat{\theta}^{(1)}) > 0.8 & \text{(Cross-validation)}
\end{cases}
\end{equation}
where $d(\cdot, \cdot)$ is distributional distance, $\epsilon_P, \epsilon_B$ are 
tolerance thresholds, and $r$ is correlation with empirical benchmarks.
\end{axiom}

\begin{axiom}[Persona Specification Protocol]
\label{ax:bbb-per}
\textbf{(BBB-PER)} Each LLM-simulated agent must be specified by a complete 
context vector:
\begin{equation}
\mathcal{P} = \langle \boldsymbol{\Psi}, \mathbf{u}_0, \mathbf{w}, \mathcal{H}, \mathcal{C} \rangle
\end{equation}
where:
\begin{itemize}
    \item $\boldsymbol{\Psi} \in [0,1]^8$ = 8-dimensional context vector
    \item $\mathbf{u}_0 \in \mathbb{R}^6$ = baseline FEPSDE utility levels
    \item $\mathbf{w} \in \Delta^5$ = utility dimension weights (simplex)
    \item $\mathcal{H}$ = decision history (past choices)
    \item $\mathcal{C}$ = constraints (budget, time, information)
\end{itemize}

\textbf{Completeness condition:}
\begin{equation}
\text{Complete}(\mathcal{P}) \Leftrightarrow \forall \psi_k \in \boldsymbol{\Psi}: \psi_k \in [0,1] \text{ is specified}
\end{equation}
\end{axiom}

\begin{axiom}[Perspective Rotation Requirement]
\label{ax:bbb-rot}
\textbf{(BBB-ROT)} To reduce systematic bias, each parameter must be estimated 
from at least $M \geq 3$ perspectives:
\begin{equation}
\hat{\theta}^{(2)} = \frac{1}{M} \sum_{m=1}^{M} \hat{\theta}^{(2)}_m \quad \text{where } M \geq 3
\end{equation}

\textbf{Required perspectives:}
\begin{enumerate}
    \item $\hat{\theta}^{(2)}_{\text{econ}}$ = Economist perspective (rational choice)
    \item $\hat{\theta}^{(2)}_{\text{psych}}$ = Psychologist perspective (behavioral)
    \item $\hat{\theta}^{(2)}_{\text{pract}}$ = Practitioner perspective (applied)
\end{enumerate}

\textbf{Consistency check:}
\begin{equation}
\text{CV}(\hat{\theta}^{(2)}_1, \ldots, \hat{\theta}^{(2)}_M) = \frac{\sigma_{\hat{\theta}}}{\bar{\theta}} < 0.3
\end{equation}
\end{axiom}

\begin{axiom}[Variation Protocol]
\label{ax:bbb-var}
\textbf{(BBB-VAR)} LLM-MC estimation requires systematic variation across:
\begin{equation}
\hat{\theta}^{(2)} = \frac{1}{|\mathcal{T}| \cdot |\mathcal{S}| \cdot N} \sum_{\tau \in \mathcal{T}} \sum_{s \in \mathcal{S}} \sum_{n=1}^{N} \hat{\theta}(\tau, s, n)
\end{equation}
where:
\begin{itemize}
    \item $\mathcal{T} = \{0.3, 0.7, 1.0\}$ = temperature settings
    \item $\mathcal{S}$ = scenario variations (at least 10)
    \item $N \geq 100$ = replications per cell
\end{itemize}

\textbf{Total minimum queries:}
\begin{equation}
N_{\text{total}} = |\mathcal{T}| \cdot |\mathcal{S}| \cdot N \cdot M \geq 3 \cdot 10 \cdot 100 \cdot 3 = 9000
\end{equation}
\end{axiom}

%------------------------------------------------------------------------------
\subsection{Tier 3 Axioms: Theoretical Constraints}
%------------------------------------------------------------------------------

\begin{axiom}[Axiomatic Constraint Priority]
\label{ax:bbb-axi}
\textbf{(BBB-AXI)} Theoretical constraints take absolute precedence:
\begin{equation}
\forall \theta \in \Theta_{\text{constrained}}: \theta = f_{\text{axiom}}(\theta) \text{ regardless of } \hat{\theta}^{(1)}, \hat{\theta}^{(2)}
\end{equation}

\textbf{Constrained parameters:}
\begin{align}
\gamma_{dd} &= 0 \quad \forall d \in \{F, E, P, S, D, Ec\} \quad \text{(No self-complementarity)} \\
\sum_d w_d &= 1 \quad \text{(Weight normalization)} \\
\delta &\in (0, 1) \quad \text{(Discount factor bounds)}
\end{align}
\end{axiom}

\begin{axiom}[Symmetry Requirement]
\label{ax:bbb-sym}
\textbf{(BBB-SYM)} The $\Gamma$-matrix and $\Lambda$-matrix must be symmetric:
\begin{equation}
\gamma_{dd'} = \gamma_{d'd} \quad \forall d, d' \in \mathcal{D}
\end{equation}
\begin{equation}
\lambda_{kl} = \lambda_{lk} \quad \forall k, l \in \{1, \ldots, 8\}
\end{equation}

\textbf{Enforcement:} When asymmetric estimates arise:
\begin{equation}
\hat{\gamma}_{dd'}^{\text{sym}} = \frac{\hat{\gamma}_{dd'} + \hat{\gamma}_{d'd}}{2}
\end{equation}

\textbf{Theoretical basis:} Young's theorem on cross-partial derivatives:
\begin{equation}
\frac{\partial^2 U}{\partial u_d \partial u_{d'}} = \frac{\partial^2 U}{\partial u_{d'} \partial u_d}
\end{equation}
\end{axiom}

\begin{axiom}[Boundedness Constraint]
\label{ax:bbb-bnd}
\textbf{(BBB-BND)} All complementarity coefficients are bounded:
\begin{equation}
\gamma_{dd'} \in [-1, 1] \quad \forall d \neq d'
\end{equation}
\begin{equation}
\lambda_{kl} \in [-1, 1] \quad \forall k \neq l
\end{equation}

\textbf{Interpretation:}
\begin{align}
\gamma = 1 &\Rightarrow \text{Perfect complementarity (Leontief)} \\
\gamma = 0 &\Rightarrow \text{Independence (additive separability)} \\
\gamma = -1 &\Rightarrow \text{Perfect substitutability}
\end{align}

\textbf{Enforcement:} Estimates outside bounds are truncated:
\begin{equation}
\hat{\gamma}^{\text{bounded}} = \max(-1, \min(1, \hat{\gamma}))
\end{equation}
\end{axiom}

%------------------------------------------------------------------------------
\subsection{Method Axioms}
%------------------------------------------------------------------------------

\begin{axiom}[Estimation Method Selection]
\label{ax:bbb-est-full}
\textbf{(BBB-EST)} Method selection follows the lexicographic ordering:
\begin{equation}
\mathcal{M}^* = \arg\max_{\mathcal{M}} \left( \text{Avail}(\mathcal{M}), \text{Causal}(\mathcal{M}), -\text{Cost}(\mathcal{M}), \text{Context}(\mathcal{M}) \right)
\end{equation}

\textbf{Method ranking:}
\begin{center}
\begin{tabular}{lcccc}
\toprule
Method & Avail & Causal & Cost & Context \\
\midrule
A (Econometric) & 2 & 4 & 3 & 2 \\
B (LLM-MC) & 5 & 2 & 5 & 4 \\
C (Experimental) & 1 & 5 & 1 & 3 \\
D (Meta-Analytic) & 3 & 3 & 4 & 2 \\
E (Milgrom) & 4 & 3 & 4 & 3 \\
\bottomrule
\end{tabular}
\end{center}
(5 = best, 1 = worst)
\end{axiom}

\begin{axiom}[Milgrom Diagnostic Validity]
\label{ax:bbb-mil-full}
\textbf{(BBB-MIL)} Milgrom diagnostic tests are valid when:
\begin{equation}
\text{Valid}_{\text{Milgrom}} \Leftrightarrow 
\begin{cases}
N_{\text{cell}} \geq 30 & \text{(Sample size per cell)} \\
\exists \theta: \text{Cov}(\theta, X) \neq 0, \text{Cov}(\theta, \varepsilon) = 0 & \text{(Exogenous variation)} \\
T \geq 3 & \text{(Panel length for Le Chatelier)}
\end{cases}
\end{equation}

\textbf{Complementarity Diagnostic Index:}
\begin{equation}
\text{CDI} = \frac{1}{6} \sum_{t=1}^{6} \mathbb{1}\{p_t < 0.05\}
\end{equation}
where $p_t$ is the p-value for test $t \in \{\text{Pairwise, Cluster, Coherent, MCS, LeCh, Momentum}\}$.
\end{axiom}

%------------------------------------------------------------------------------
\subsection{Validation Axioms}
%------------------------------------------------------------------------------

\begin{axiom}[Cross-Validation Protocol]
\label{ax:bbb-val}
\textbf{(BBB-VAL)} Every parameter estimate must satisfy cross-validation:
\begin{equation}
R^2_{\text{OOS}} = 1 - \frac{\sum_{i \in \text{test}} (Y_i - \hat{Y}_i(\hat{\theta}))^2}{\sum_{i \in \text{test}} (Y_i - \bar{Y})^2} \geq R^2_{\min}
\end{equation}
where $R^2_{\min} = 0.5$ for Tier 1, $R^2_{\min} = 0.4$ for Tier 2.

\textbf{K-fold protocol:}
\begin{equation}
\bar{R}^2_{\text{CV}} = \frac{1}{K} \sum_{k=1}^{K} R^2_k \quad \text{with } K = 5
\end{equation}

\textbf{Stability criterion:}
\begin{equation}
\text{CV}(R^2_1, \ldots, R^2_K) < 0.2
\end{equation}
\end{axiom}

\begin{axiom}[Robustness Requirements]
\label{ax:bbb-rob}
\textbf{(BBB-ROB)} Estimates must be robust to specification changes:
\begin{equation}
\text{Robust}(\hat{\theta}) \Leftrightarrow \forall s \in \mathcal{S}_{\text{spec}}: 
\frac{|\hat{\theta}_s - \hat{\theta}_{\text{base}}|}{\sigma_{\hat{\theta}}} < 2
\end{equation}

\textbf{Required robustness checks:}
\begin{enumerate}
    \item Alternative control variables: $|\Delta \hat{\theta}| < 2\sigma$
    \item Sample restrictions (drop outliers): $|\Delta \hat{\theta}| < 2\sigma$
    \item Functional form (linear vs. log): $|\Delta \hat{\theta}| < 2\sigma$
    \item Time period variation: $|\Delta \hat{\theta}| < 2\sigma$
    \item Clustering adjustment: $|\Delta \hat{\theta}| < 2\sigma$
\end{enumerate}

\textbf{Robustness score:}
\begin{equation}
\text{RS}(\hat{\theta}) = \frac{\#\{\text{checks passed}\}}{5} \geq 0.8
\end{equation}
\end{axiom}

%==============================================================================
% AXIOM INDEX
%==============================================================================

\subsection{Complete Axiom Index}
\label{sec:bbb-axiom-index}

\begin{longtable}{llp{7cm}l}
\toprule
\textbf{Code} & \textbf{Name} & \textbf{Core Statement} & \textbf{Section} \\
\midrule
\endfirsthead
\toprule
\textbf{Code} & \textbf{Name} & \textbf{Core Statement} & \textbf{Section} \\
\midrule
\endhead
\multicolumn{4}{l}{\textbf{Epistemological Axioms}} \\
BBB-EPI & Epistemic Tag & $E(\theta) = \sum w_s q_s r_s / \sum w_s$ & \ref{ax:bbb-epi} \\
BBB-HIE & Hierarchy Precedence & Tier 1 $>$ Tier 2 $>$ Tier 3 & \ref{ax:bbb-hie} \\
BBB-TRA & Transparency & Prov($\theta$) = $\langle \hat{\theta}, \sigma, E, \mathcal{S}, t, \mathcal{M} \rangle$ & \ref{ax:bbb-tra} \\
BBB-UPD & Update Obligation & Update when $D_{KL} > \kappa$ & \ref{ax:bbb-upd} \\
\midrule
\multicolumn{4}{l}{\textbf{Tier 1 Axioms (Empirical)}} \\
BBB-EMP & Empirical Validity & $N \geq N_{\min}$, Var$(X) > 0$, $\mathbb{E}[\varepsilon|X] = 0$ & \ref{ax:bbb-emp} \\
BBB-CAU & Causal Identification & IV relevance + exclusion + independence & \ref{ax:bbb-cau} \\
BBB-REP & Replicability & $|\hat{\theta}_{\text{orig}} - \hat{\theta}_{\text{rep}}| / \sigma < 1.96$ & \ref{ax:bbb-rep} \\
\midrule
\multicolumn{4}{l}{\textbf{Tier 2 Axioms (LLM-MC)}} \\
BBB-LLM & LLM Validity & Persona match + unbiased + $r > 0.8$ & \ref{ax:bbb-llm-full} \\
BBB-PER & Persona Specification & $\mathcal{P} = \langle \boldsymbol{\Psi}, \mathbf{u}_0, \mathbf{w}, \mathcal{H}, \mathcal{C} \rangle$ & \ref{ax:bbb-per} \\
BBB-ROT & Perspective Rotation & $M \geq 3$ perspectives, CV $< 0.3$ & \ref{ax:bbb-rot} \\
BBB-VAR & Variation Protocol & $N_{\text{total}} \geq 9000$ queries & \ref{ax:bbb-var} \\
\midrule
\multicolumn{4}{l}{\textbf{Tier 3 Axioms (Theoretical)}} \\
BBB-AXI & Axiomatic Priority & Constraints override empirical estimates & \ref{ax:bbb-axi} \\
BBB-SYM & Symmetry & $\gamma_{dd'} = \gamma_{d'd}$, $\lambda_{kl} = \lambda_{lk}$ & \ref{ax:bbb-sym} \\
BBB-BND & Boundedness & $\gamma, \lambda \in [-1, 1]$ & \ref{ax:bbb-bnd} \\
\midrule
\multicolumn{4}{l}{\textbf{Method Axioms}} \\
BBB-EST & Method Selection & Lexicographic: Avail $>$ Causal $>$ Cost $>$ Context & \ref{ax:bbb-est-full} \\
BBB-MIL & Milgrom Validity & $N \geq 30$/cell, exogenous $\theta$, $T \geq 3$ & \ref{ax:bbb-mil-full} \\
\midrule
\multicolumn{4}{l}{\textbf{Validation Axioms}} \\
BBB-VAL & Cross-Validation & $R^2_{\text{OOS}} \geq 0.5$ (Tier 1), $\geq 0.4$ (Tier 2) & \ref{ax:bbb-val} \\
BBB-ROB & Robustness & $|\Delta\hat{\theta}| < 2\sigma$ across 5 checks & \ref{ax:bbb-rob} \\
\bottomrule
\caption{Complete axiom index for Appendix EST}
\label{tab:bbb-axiom-index}
\end{longtable}




%==============================================================================
% LIMITATIONS AND CAVEATS
%==============================================================================

\section{Limitations and Caveats}
\label{sec:bbb-limitations}

Transparent acknowledgment of methodological limitations is essential for 
responsible use of EBF parameter estimates.

\subsection{Tier 1 Limitations: Empirical Estimation}

\begin{enumerate}
    \item \textbf{External Validity:}
    Estimates from specific populations (e.g., Swiss, US) may not generalize 
    to other contexts. The $\delta$-tensor attempts to capture this heterogeneity, 
    but unobserved context factors may remain.
    \begin{equation}
    \text{Generalization Error} = |\gamma_{\text{source}} - \gamma_{\text{target}}| \cdot (1 - \text{Context Overlap})
    \end{equation}
    
    \item \textbf{Identification Challenges:}
    True complementarity requires exogenous variation in \textit{both} dimensions 
    simultaneously. Many ``complementarity'' estimates may reflect omitted variable 
    bias or reverse causality.
    
    \item \textbf{Measurement Error:}
    Utility dimensions (especially Emotional, Social) are measured with error. 
    Attenuation bias implies:
    \begin{equation}
    \text{plim } \hat{\gamma} = \gamma \cdot \frac{\sigma^2_X}{\sigma^2_X + \sigma^2_\varepsilon} < \gamma
    \end{equation}
    
    \item \textbf{Temporal Instability:}
    Complementarity patterns may shift over time due to technological change, 
    cultural evolution, or policy shifts. Estimates should be re-validated 
    periodically (recommended: every 3-5 years).
\end{enumerate}

\subsection{Tier 2 Limitations: LLM Monte Carlo}

\begin{enumerate}
    \item \textbf{Hallucination Risk:}
    LLMs may generate plausible but factually incorrect behavioral predictions. 
    Mitigation: Cross-validation with empirical benchmarks (Axiom BBB-LLM).
    \begin{equation}
    P(\text{Hallucination}) \approx 1 - r(\hat{\theta}^{(2)}, \hat{\theta}^{(1)})
    \end{equation}
    
    \item \textbf{Training Data Bias:}
    LLM priors reflect the corpus they were trained on, which overrepresents 
    WEIRD (Western, Educated, Industrialized, Rich, Democratic) populations.
    
    \item \textbf{Prompt Sensitivity:}
    Estimates can vary significantly with prompt phrasing. Mitigation: 
    Perspective rotation (Axiom BBB-ROT) and scenario variation (Axiom BBB-VAR).
    
    \item \textbf{Reproducibility Concerns:}
    LLM outputs may change with model updates. All estimates should record:
    \begin{itemize}
        \item Model version (e.g., Claude 3.5, GPT-4)
        \item Timestamp
        \item Exact prompts used
    \end{itemize}
    
    \item \textbf{Ceiling on Epistemic Status:}
    By construction, LLM-MC estimates cannot exceed $E(\theta) = 0.9$. 
    They are \textit{priors}, not ground truth.
\end{enumerate}

\subsection{Tier 3 Limitations: Theoretical Constraints}

\begin{enumerate}
    \item \textbf{Over-Constraint Risk:}
    Forcing parameters to exact values (e.g., $\gamma_{dd} = 0$) may miss 
    real-world deviations. Theory is a map, not the territory.
    
    \item \textbf{Model Misspecification:}
    If the underlying utility function is misspecified, axiomatic constraints 
    may propagate errors throughout the system.
\end{enumerate}

\subsection{General Limitations}

\begin{enumerate}
    \item \textbf{Dimensionality:}
    With 185 parameters, overfitting is a risk even with large samples. 
    Regularization and out-of-sample validation are essential.
    
    \item \textbf{Context Drift:}
    The $\Psi$-dimensions themselves may shift meaning over time. 
    ``Digital Access'' meant something different in 2010 vs. 2025.
    
    \item \textbf{Aggregation Paradoxes:}
    Individual-level complementarities may not hold at aggregate levels 
    due to composition effects (Simpson's paradox).
    
    \item \textbf{Interaction Complexity:}
    The model assumes pairwise interactions ($\gamma_{dd'}$). Higher-order 
    interactions ($\gamma_{dd'd''}$) are not captured but may matter.
\end{enumerate}

\subsection{Recommendations for Users}

\begin{tcolorbox}[colback=yellow!5!white, colframe=yellow!75!black, title=Best Practices]
\begin{enumerate}
    \item \textbf{Check Epistemic Status:} Always examine $E(\theta)$ before using estimates.
    \item \textbf{Use Confidence Intervals:} Point estimates alone are insufficient.
    \item \textbf{Validate Locally:} Run pilots to validate parameters in your specific context.
    \item \textbf{Update Regularly:} Treat estimates as priors, not constants.
    \item \textbf{Report Uncertainty:} Communicate estimation uncertainty to stakeholders.
    \item \textbf{Triangulate:} Use multiple methods where possible.
\end{enumerate}
\end{tcolorbox}

%==============================================================================
% EXPANDED CROSS-REFERENCES
%==============================================================================

\section{Cross-References to EBF Documentation}
\label{sec:bbb-crossref-expanded}

Appendix EST serves as the \textbf{methodological hub} of the EBF framework, 
connecting to all other core documentation.

\subsection{Upstream Dependencies (BBB receives from)}

\begin{longtable}{lp{4cm}p{6cm}}
\toprule
\textbf{Source} & \textbf{What BBB Receives} & \textbf{How BBB Uses It} \\
\midrule
\endfirsthead
Appendix HOW (\ref{app:complementarity-how}) & Complementarity theory, $\gamma$ definitions & Defines \textit{what} parameters mean \\
Appendix CTW (\ref{app:psi-dimensions}) & 8$\Psi$-dimension definitions & Context vector for $\delta$-modulation \\
Appendix WAT (\ref{app:fepsde-matrix}) & FEPSDE utility structure & Defines dimensions $d \in \mathcal{D}$ \\
Appendix AWA (\ref{app:axiom-formalization}) & Axiom formalization & Tier 3 axiomatic constraints \\
Appendix LLM1 (\ref{app:llm-monte-carlo}) & LLM-MC protocol details & Method B implementation \\
Appendix OPS (\ref{app:operationalization}) & Empirical data sources & Tier 1 calibration inputs \\
Chapter 4 & Application cases & Validation benchmarks \\
\bottomrule
\caption{Upstream dependencies for Appendix EST}
\end{longtable}

\subsection{Downstream Outputs (BBB provides to)}

\begin{longtable}{lp{4cm}p{6cm}}
\toprule
\textbf{Target} & \textbf{What BBB Provides} & \textbf{How Target Uses It} \\
\midrule
\endfirsthead
Appendix HOW (\ref{app:complementarity-how}) & Calibrated $\Gamma$, $\Lambda$, $\delta$ & Quantified complementarity for analysis \\
Appendix WHO (\ref{app:aggregation-levels}) & Level-specific parameters & Parameters by aggregation level \\
Appendix WAT (\ref{app:fepsde-matrix}) & Calibrated weights $w_d$ & FEPSDE utility specification \\
Chapter 3 & Parameter values & Model implementation \\
Chapter 5 & Intervention effect sizes & Policy recommendations \\
BEATRIX System & Complete $\Theta$-vector & Computational implementation \\
\bottomrule
\caption{Downstream outputs from Appendix EST}
\end{longtable}

\subsection{Bidirectional Cross-Reference Matrix}

\begin{center}
\small
\begin{tabular}{l|ccccccc}
\toprule
& \textbf{B} & \textbf{V} & \textbf{C} & \textbf{AAA} & \textbf{AN} & \textbf{E} & \textbf{AU} \\
\midrule
\textbf{BBB receives} & Theory & $\Psi$ & FEPSDE & Levels & LLM-MC & Data & Axioms \\
\textbf{BBB provides} & $\Gamma,\Lambda,\delta$ & -- & $w_d$ & Params & Validation & -- & -- \\
\midrule
\textbf{Link strength} & Strong & Strong & Medium & Medium & Strong & Medium & Weak \\
\bottomrule
\end{tabular}
\end{center}

\subsection{The EBF Workflow}

\begin{center}
\fbox{\parbox{0.9\textwidth}{
\centering
\textbf{EBF Parameter Workflow}\\[1em]
\texttt{Context (V)} $\xrightarrow{\Psi}$ 
\texttt{\textbf{BBB}} $\xrightarrow{\Gamma,\Lambda,\delta}$ 
\texttt{Analysis (B)} $\xrightarrow{\text{effects}}$ 
\texttt{Intervention}\\[0.5em]
$\uparrow$\hspace{8cm}$\downarrow$\\
\texttt{Update (BBB-UPD)} $\xleftarrow{\text{outcomes}}$ \texttt{Feedback}
}}
\end{center}

\textbf{Workflow Description:}
\begin{enumerate}
    \item \textbf{Context Assessment:} Measure $\boldsymbol{\Psi}$ using Appendix CTW (\ref{app:psi-dimensions}) instruments
    \item \textbf{Parameter Lookup:} Retrieve $\Gamma$, $\Lambda$ from BBB Parameter Dictionary
    \item \textbf{Context Adjustment:} Compute $\gamma(\boldsymbol{\Psi})$ using $\delta$-tensor
    \item \textbf{Complementarity Analysis:} Apply Appendix HOW (\ref{app:complementarity-how}) theory to identify leverage points
    \item \textbf{Intervention Design:} Use Chapter 5 protocols
    \item \textbf{Feedback Loop:} Update BBB estimates with intervention outcomes (Axiom BBB-UPD)
\end{enumerate}

\subsection{Quick Reference: Where to Find What}

\begin{center}
\begin{tabular}{lll}
\toprule
\textbf{If you need...} & \textbf{Go to...} & \textbf{Section} \\
\midrule
Theory of complementarity & Appendix HOW (\ref{app:complementarity-how}) & Teil I-II \\
$\gamma_{dd'}$ values & \textbf{BBB} & Parameter Dictionary \\
$\lambda_{kl}$ values & \textbf{BBB} & Parameter Dictionary \\
$\delta_{dd'k}$ values & \textbf{BBB} & $\delta$-Tensor \\
$\Psi$-dimension definitions & Appendix CTW (\ref{app:psi-dimensions}) & All \\
FEPSDE definitions & Appendix WAT (\ref{app:fepsde-matrix}) & Teil I \\
Estimation method details & \textbf{BBB} & Five Methods \\
LLM-MC protocol & Appendix LLM1 (\ref{app:llm-monte-carlo}) & Protocol \\
Axiom statements & \textbf{BBB} + AU & Axiom System \\
Worked examples & \textbf{BBB} + F & Worked Example \\
\bottomrule
\end{tabular}
\end{center}



\section{Glossary of Key Terms}
\label{sec:bbb-glossary}

\begin{description}

\item[Bayesian Updating]
The process of revising probability estimates for parameters as new evidence 
arrives: $P(\theta|\text{Data}) \propto P(\text{Data}|\theta) \cdot P(\theta)$.
In EBF, Tier 2 estimates serve as priors that are updated with Tier 1 data.

\item[Calibration]
The process of determining parameter values by matching model predictions to 
observed data. Distinct from estimation in that it may involve judgment and 
iterative adjustment rather than pure statistical inference.

\item[Complementarity Diagnostic Index (CDI)]
Aggregate score from 0 to 1 indicating strength of evidence for complementarity 
based on six Milgrom diagnostic tests. CDI $\geq 0.83$ indicates strong evidence.

\item[Confidence Score]
Measure of certainty in a predicted intervention effect: 
$\text{Confidence}(\tau) = 1 - \sigma_\tau/|\hat{\tau}|$. 
Used to determine whether to scale, pilot, or research further.

\item[Context Modulation]
The phenomenon whereby complementarity coefficients $\gamma_{dd'}$ vary with 
context $\boldsymbol{\Psi}$: $\gamma_{dd'}(\boldsymbol{\Psi}) = \bar{\gamma}_{dd'} + \sum_k \delta_{dd'k} \Psi_k$.

\item[Econometric Estimation]
Method A: Estimating parameters from panel data using interaction regressions 
with fixed effects. Gold standard for causal identification when data exists.

\item[Epistemic Tag $E(\theta)$]
A value in $[0, 1]$ assigned to each parameter indicating the reliability of 
its source: $E \approx 1.0$ for direct measurement, $E \approx 0.7-0.9$ for 
LLM-MC, $E < 0.5$ for theoretical constraints.

\item[Epistemic Uncertainty]
Uncertainty arising from incomplete knowledge (reducible with more data), 
as opposed to aleatory uncertainty (inherent randomness). Represented by 
$\sigma^2_\theta$ in the parameter distribution.

\item[Experimental Estimation]
Method C: Using randomized factorial designs to identify complementarity via 
the interaction effect $\gamma = [Y_{11} - Y_{10}] - [Y_{01} - Y_{00}]$.

\item[Factorial Design]
Experimental design where multiple treatments are varied independently, 
allowing identification of interaction effects (complementarities) from 
the difference-in-differences of treatment effects.

\item[FEPSDE]
The six utility dimensions in EBF: Financial, Emotional, Physical, 
Social, Digital, Ecological. Defined in Appendix WAT (\ref{app:fepsde-matrix}).

\item[$\Gamma$-Matrix]
The $6 \times 6$ symmetric matrix of utility-utility complementarity 
coefficients $[\gamma_{dd'}]$ for the FEPSDE dimensions. 15 unique parameters.

\item[Hard Empirical (Tier 1)]
The highest reliability tier of evidence: direct econometric estimates from 
experimental data, meta-analyses, or high-quality surveys. $E(\theta) \approx 0.9-1.0$.

\item[Hierarchy of Evidence]
The three-tier system ranking parameter sources by reliability: 
Tier 1 (Empirical) $>$ Tier 2 (LLM-MC) $>$ Tier 3 (Axiomatic).

\item[Homo Silicus]
Term coined by Horton (2023) for LLM-simulated economic agents used in 
synthetic behavioral research.

\item[Interaction Regression]
Regression model including product terms $X_d \cdot X_{d'}$ to estimate 
complementarity coefficients: $Y = \alpha + \sum_d \beta_d X_d + \sum_{d<d'} \gamma_{dd'} X_d X_{d'} + \varepsilon$.

\item[$\Lambda$-Matrix]
The $8 \times 8$ symmetric matrix of context-context complementarity 
coefficients $[\lambda_{kl}]$ for the 8$\Psi$-dimensions. 28 unique parameters.

\item[Le Chatelier Test]
Milgrom diagnostic comparing short-run and long-run responses to a shock. 
Under complementarity, $|\beta^{LR}| > |\beta^{SR}|$.

\item[LLM Monte Carlo (LLM-MC)]
Method B: Using Large Language Models to generate synthetic behavioral data 
through persona simulation and scenario variation. Tier 2 evidence source.

\item[Meta-Analytic Calibration]
Method D: Bayesian synthesis of multiple studies to derive parameter estimates, 
starting with theoretical priors and updating with each study's evidence.

\item[Milgrom Diagnostic Tests]
Method E: Six indirect tests for complementarity derived from 
Brynjolfsson \& Milgrom (2013): Pairwise, Clustering, Coherent Direction, 
Monotone CS, Le Chatelier, Momentum.

\item[Monotone Comparative Statics (MCS)]
Technique for deriving qualitative predictions about how optimal choices 
change with parameters, without solving the optimization problem explicitly.

\item[Parameter Dictionary]
The complete inventory of EBF parameters with values, standard errors, 
source tiers, and epistemic status. Central reference in Appendix EST.

\item[Parameter Universe ($\Theta$)]
The complete set of 185 parameters required by EBF: 15 ($\Gamma$) + 28 ($\Lambda$) 
+ 120 ($\delta$) + 6 ($w$) + 6 ($r$) + 1 ($\delta$) + 8 ($\Psi^*$) + 1 ($K^*$).

\item[Perspective Rotation]
LLM-MC technique of querying the same scenario from multiple expert viewpoints 
(economist, psychologist, practitioner) to reduce bias and increase robustness.

\item[Prior]
In Bayesian inference, the probability distribution representing beliefs about 
a parameter before observing data. In EBF, Tier 2/3 estimates serve as priors.

\item[Posterior]
The updated probability distribution for a parameter after incorporating new 
evidence via Bayes' rule.

\item[$\Psi$-Dimensions]
The eight context dimensions in EBF: Institutional ($\Psi_I$), Social ($\Psi_S$), 
Cognitive ($\Psi_C$), Informational ($\Psi_K$), Economic ($\Psi_E$), Temporal ($\Psi_T$), 
Market Scope ($\Psi_M$), Factor Flexibility ($\Psi_F$). Defined in Appendix CTW (\ref{app:psi-dimensions}).

\item[Supermodularity]
Mathematical property where $f(x \vee y) + f(x \wedge y) \geq f(x) + f(y)$, 
implying that increasing one variable raises the marginal return to others. 
The formal basis for complementarity.

\item[Synthetic Priors (Tier 2)]
Parameter estimates derived from LLM Monte Carlo, representing compressed 
scientific knowledge. Medium-high reliability: $E(\theta) \approx 0.7-0.9$.

\item[Theoretical Constraints (Tier 3)]
Parameter values determined by logical necessity (e.g., symmetry, budget 
constraints, identity properties). Absolute reliability for the constraint 
itself, but low granularity.

\item[Three-Tier Strategy]
EBF's hybrid approach to parameter estimation combining empirical data 
(Tier 1), LLM-MC synthesis (Tier 2), and axiomatic constraints (Tier 3).

\item[Tier Integration Protocol]
The decision procedure for selecting parameter sources: check Tier 1 first, 
then Tier 3 constraints, default to Tier 2, document all choices.

\item[Treatment Effect Heterogeneity]
Variation in intervention effects across contexts: $\tau(\boldsymbol{\Psi})$ 
rather than a single $\tau$. The central motivation for context-dependent modeling.

\item[Trilemma]
The fundamental challenge in behavioral modeling: theoretical richness vs. 
data scarcity vs. quantification need. EBF resolves it via the three-tier strategy.

\item[Validation]
The process of confirming that estimated parameters produce accurate predictions 
on held-out data and are consistent across estimation methods.

\end{description}




%==============================================================================
% PART 4: BACK MATTER (Ergänzungen)
%==============================================================================

%------------------------------------------------------------------------------
% 4.1 SUMMARY BOX
%------------------------------------------------------------------------------

\section{Summary}
\label{sec:bbb-summary}

\begin{tcolorbox}[colback=blue!5!white, colframe=blue!75!black, 
    title=\textbf{Appendix EST Summary: The Epistemics of Numbers}]

\textbf{CORE Question Answered:} WHERE do the numbers come from?

\textbf{The Fundamental Trilemma:}
\begin{center}
\textit{Theoretical richness} $\leftrightarrow$ \textit{Data scarcity} $\leftrightarrow$ \textit{Quantification need}
\end{center}

\textbf{The Solution: Three-Tier Hybrid Strategy}
\begin{enumerate}[nosep]
    \item \textbf{Tier 1} (Hard Empirical): Meta-analyses, RCTs, panel data $\rightarrow E(\theta) \approx 0.9$
    \item \textbf{Tier 2} (LLM Monte Carlo): Synthetic priors from compressed knowledge $\rightarrow E(\theta) \approx 0.8$
    \item \textbf{Tier 3} (Axiomatic): Theoretical constraints (symmetry, bounds) $\rightarrow E(\theta) = 1.0$
\end{enumerate}

\textbf{Five Estimation Methods:}
\begin{center}
\begin{tabular}{@{}clcl@{}}
A & Econometric & B & LLM Monte Carlo \\
C & Experimental & D & Meta-Analytic \\
E & Milgrom Diagnostic & & \\
\end{tabular}
\end{center}

\textbf{Key Outputs:}
\begin{itemize}[nosep]
    \item $\Gamma$-matrix: 15 utility-utility complementarity coefficients
    \item $\Lambda$-matrix: 28 context-context complementarity coefficients
    \item $\delta$-tensor: 120 context modulation coefficients
    \item 18 axioms formalizing estimation validity
    \item Complete provenance for every parameter
\end{itemize}

\textbf{Core Principle:} \textit{``Every number has an address''} --- transparent, 
updatable, epistemically honest parameter infrastructure.

\end{tcolorbox}

%------------------------------------------------------------------------------
% 4.2 CRITICAL FOUNDATIONS (für CORE Appendices)
%------------------------------------------------------------------------------

\section{Critical Foundations}
\label{sec:bbb-critical-foundations}

\begin{tcolorbox}[colback=red!5!white, colframe=red!75!black, 
    title=\textbf{Critical Foundations for EBF}]

The following elements from Appendix EST are \textbf{foundational} for the 
entire EBF framework. Changes to these require careful propagation through 
all dependent components.

\textbf{Foundation 1: The Epistemic Tag System}
\begin{equation}
E(\theta) \in [0, 1] \text{ for all } \theta \in \Theta
\end{equation}
Every parameter must have an epistemic status. This enables principled 
uncertainty quantification throughout the model.

\textbf{Foundation 2: The Three-Tier Hierarchy}
\begin{equation}
\text{Tier 1} > \text{Tier 2} > \text{Tier 3} \quad \text{(strict precedence)}
\end{equation}
When estimates conflict, higher tiers always dominate. This prevents 
cherry-picking and ensures scientific integrity.

\textbf{Foundation 3: The Provenance Requirement}
\begin{equation}
\forall \theta \in \Theta: \text{Prov}(\theta) \text{ must be documented}
\end{equation}
No parameter may be used without complete provenance (value, SE, source, 
method, timestamp). This enables reproducibility and updating.

\textbf{Foundation 4: The Context Modulation Structure}
\begin{equation}
\gamma_{dd'}(\boldsymbol{\Psi}) = \bar{\gamma}_{dd'} + \sum_k \delta_{dd'k} (\Psi_k - \bar{\Psi}_k)
\end{equation}
Complementarity is not constant but varies systematically with context. 
The $\delta$-tensor is essential for cross-context generalization.

\textbf{Foundation 5: Bayesian Updating Commitment}
\begin{equation}
\theta^{(t+1)} = f(\theta^{(t)}, D_{\text{new}}) \text{ when } I(D_{\text{new}}; \theta) > \kappa
\end{equation}
Parameters are priors, not constants. The framework commits to continuous 
updating as new evidence arrives.

\end{tcolorbox}

%------------------------------------------------------------------------------
% 4.3 OPEN ISSUES / FUTURE WORK
%------------------------------------------------------------------------------

\section{Open Issues and Future Work}
\label{sec:bbb-open-issues}

\subsection{Methodological Open Questions}

\begin{enumerate}
    \item \textbf{LLM Model Dependence:}
    Current Tier 2 estimates are based on specific LLM versions (Claude 3.5, GPT-4).
    How stable are estimates across model generations? Systematic comparison needed.
    
    \item \textbf{Higher-Order Interactions:}
    The current framework captures pairwise complementarities ($\gamma_{dd'}$) but not 
    three-way interactions ($\gamma_{dd'd''}$). Extension to higher-order terms may 
    improve fit but at cost of parsimony.
    
    \item \textbf{Dynamic Parameter Evolution:}
    How do complementarity coefficients change over time? Need longitudinal 
    re-estimation protocol (currently: every 3-5 years recommended).
    
    \item \textbf{Extreme Context Extrapolation:}
    $\delta$-modulation is validated for $\Psi \in [0.2, 0.8]$. Behavior at 
    extremes ($\Psi < 0.1$ or $\Psi > 0.9$) may be non-linear.
\end{enumerate}

\subsection{Data Collection Priorities}

\begin{enumerate}
    \item \textbf{Tier 1 Expansion:} Currently only 18/120 $\delta$-coefficients 
    have Tier 1 estimates. Priority: cross-country experimental panels.
    
    \item \textbf{Non-WEIRD Populations:} LLM training data overrepresents 
    Western populations. Need targeted LLM-MC calibration for Africa, South Asia.
    
    \item \textbf{Ecological Dimension:} $\gamma_{*Ec}$ estimates have lowest 
    epistemic status. Climate economics data integration needed.
\end{enumerate}

\subsection{Technical Extensions}

\begin{enumerate}
    \item \textbf{Automated Updating Pipeline:} Build infrastructure for 
    continuous Bayesian updating as new studies are published.
    
    \item \textbf{Sensitivity Analysis Toolkit:} Software for propagating 
    parameter uncertainty through intervention predictions.
    
    \item \textbf{Cross-Validation Database:} Systematic record of LLM-MC 
    vs. empirical comparisons for ongoing validation.
\end{enumerate}

%------------------------------------------------------------------------------
% 4.4 IMPLICATIONS FOR PRACTICE
%------------------------------------------------------------------------------

\section{Implications for Practice}
\label{sec:bbb-implications}

\begin{tcolorbox}[colback=yellow!5!white, colframe=yellow!75!black, 
    title=\textbf{For Practitioners: Using BBB Estimates}]

\textbf{DO:}
\begin{itemize}[nosep]
    \item Always check $E(\theta)$ before using a parameter
    \item Use confidence intervals, not just point estimates
    \item Validate key parameters with local pilot data
    \item Document which parameter values you used (version, date)
    \item Run sensitivity analyses for parameters with $E < 0.8$
\end{itemize}

\textbf{DON'T:}
\begin{itemize}[nosep]
    \item Treat Tier 2 estimates as ground truth
    \item Extrapolate far outside the $\Psi$-range of your context
    \item Ignore the $\delta$-adjustment for your specific context
    \item Use parameters without checking for recent updates
\end{itemize}

\textbf{Decision Thresholds:}
\begin{center}
\begin{tabular}{@{}lcl@{}}
\toprule
\textbf{Confidence} & \textbf{Range} & \textbf{Action} \\
\midrule
High & $> 0.8$ & Scale intervention \\
Medium & $0.5 - 0.8$ & Run pilot first \\
Low & $< 0.5$ & More research needed \\
\bottomrule
\end{tabular}
\end{center}

\end{tcolorbox}




%------------------------------------------------------------------------------
% CENTRAL REFERENCES (Required per Template)
%------------------------------------------------------------------------------

\begin{tcolorbox}[colback=blue!5!white, colframe=blue!75!black, 
    title=\textbf{Central References}]

\textbf{Glossary:} For complete symbol definitions across the EBF framework, 
see \textbf{Appendix GLS} (\textit{Glossary and Notation}).

\textbf{Bibliography:} For full reference details, see the 
\textbf{Master Bibliography} (\texttt{bcm2\_master\_references.bib}) 
in the repository.

All symbols used in this appendix are consistent with Appendix GLS definitions.

\end{tcolorbox}

\section{References}
\label{sec:bbb-references}

\subsection{Core Methodological Sources}

\begin{itemize}

\item \textbf{Brynjolfsson, E. \& Milgrom, P. (2013).} 
``Complementarity in Organizations.'' 
In R. Gibbons \& J. Roberts (Eds.), \textit{Handbook of Organizational Economics}. 
Princeton University Press.
\textit{[Ten theorems for detecting and measuring complementarity]}

\item \textbf{Milgrom, P. \& Roberts, J. (1990).} 
``The Economics of Modern Manufacturing: Technology, Strategy, and Organization.'' 
\textit{American Economic Review}, 80(3), 511-528.
\textit{[Foundational theory of complementarities in production]}

\item \textbf{Milgrom, P. \& Shannon, C. (1994).} 
``Monotone Comparative Statics.'' 
\textit{Econometrica}, 62(1), 157-180.
\textit{[Mathematical foundations for MCS without calculus]}

\item \textbf{Topkis, D. M. (1998).} 
\textit{Supermodularity and Complementarity}. 
Princeton University Press.
\textit{[Definitive mathematical treatment of lattice theory in economics]}

\end{itemize}

\subsection{Econometric Estimation (Method A)}

\begin{itemize}

\item \textbf{Ichniowski, C., Shaw, K. \& Prennushi, G. (1997).} 
``The Effects of Human Resource Management Practices on Productivity.'' 
\textit{American Economic Review}, 87(3), 291-313.
\textit{[Steel finishing lines: bundled HR practices]}

\item \textbf{Bloom, N. \& Van Reenen, J. (2007).} 
``Measuring and Explaining Management Practices Across Firms and Countries.'' 
\textit{Quarterly Journal of Economics}, 122(4), 1351-1408.
\textit{[Cross-country management complementarity]}

\item \textbf{Bresnahan, T., Brynjolfsson, E. \& Hitt, L. (2002).} 
``Information Technology, Workplace Organization, and the Demand for Skilled Labor.'' 
\textit{Quarterly Journal of Economics}, 117(1), 339-376.
\textit{[IT-skills-organization three-way complementarity]}

\item \textbf{Athey, S. \& Stern, S. (1998).} 
``An Empirical Framework for Testing Theories about Complementarity in Organizational Design.'' 
NBER Working Paper 6600.
\textit{[Econometric identification of complementarities]}

\end{itemize}

\subsection{LLM Monte Carlo (Method B)}

\begin{itemize}

\item \textbf{Horton, J. J. (2023).} 
``Large Language Models as Simulated Economic Agents: What Can We Learn from Homo Silicus?'' 
NBER Working Paper 31122.
\textit{[LLMs as behavioral simulators]}

\item \textbf{Argyle, L. P. et al. (2023).} 
``Out of One, Many: Using Language Models to Simulate Human Samples.'' 
\textit{Political Analysis}, 31(3), 337-351.
\textit{[Perspective rotation methodology]}

\item \textbf{Aher, G., Arriaga, R. I. \& Kalai, A. T. (2023).} 
``Using Large Language Models to Simulate Multiple Humans and Replicate Human Subject Studies.'' 
\textit{Proceedings of ICML 2023}.
\textit{[LLM replication of classic experiments]}

\item \textbf{Brand, J., Israeli, A. \& Ngwe, D. (2023).} 
``Using GPT for Market Research.'' 
Harvard Business School Working Paper 23-062.
\textit{[Synthetic consumer research]}

\end{itemize}

\subsection{Experimental Methods (Method C)}

\begin{itemize}

\item \textbf{Banerjee, A. \& Duflo, E. (2009).} 
``The Experimental Approach to Development Economics.'' 
\textit{Annual Review of Economics}, 1, 151-178.
\textit{[Factorial designs in field experiments]}

\item \textbf{Muralidharan, K. \& Niehaus, P. (2017).} 
``Experimentation at Scale.'' 
\textit{Journal of Economic Perspectives}, 31(4), 103-124.
\textit{[Large-scale factorial RCTs]}

\item \textbf{Imbens, G. W. \& Rubin, D. B. (2015).} 
\textit{Causal Inference for Statistics, Social, and Biomedical Sciences}. 
Cambridge University Press.
\textit{[Factorial design identification]}

\end{itemize}

\subsection{Meta-Analytic Methods (Method D)}

\begin{itemize}

\item \textbf{Gelman, A. et al. (2013).} 
\textit{Bayesian Data Analysis} (3rd ed.). 
Chapman \& Hall/CRC.
\textit{[Bayesian updating methodology]}

\item \textbf{Vivalt, E. (2020).} 
``How Much Can We Generalize from Impact Evaluations?'' 
\textit{Journal of the European Economic Association}, 18(6), 3045-3089.
\textit{[Treatment effect heterogeneity in meta-analysis]}

\item \textbf{Meager, R. (2019).} 
``Understanding the Average Impact of Microcredit Expansions.'' 
\textit{American Economic Journal: Applied Economics}, 11(1), 57-91.
\textit{[Bayesian hierarchical meta-analysis]}

\end{itemize}

\subsection{Behavioral Economics Foundations}

\begin{itemize}

\item \textbf{Kahneman, D. \& Tversky, A. (1979).} 
``Prospect Theory: An Analysis of Decision under Risk.'' 
\textit{Econometrica}, 47(2), 263-291.
\textit{[Reference dependence, loss aversion]}

\item \textbf{Thaler, R. H. \& Sunstein, C. R. (2008).} 
\textit{Nudge: Improving Decisions about Health, Wealth, and Happiness}. 
Yale University Press.
\textit{[Choice architecture, behavioral interventions]}

\item \textbf{Fehr, E. \& Schmidt, K. M. (1999).} 
``A Theory of Fairness, Competition, and Cooperation.'' 
\textit{Quarterly Journal of Economics}, 114(3), 817-868.
\textit{[Social preferences, inequity aversion]}

\item \textbf{Loewenstein, G. (1996).} 
``Out of Control: Visceral Influences on Behavior.'' 
\textit{Organizational Behavior and Human Decision Processes}, 65(3), 272-292.
\textit{[Visceral factors, hot-cold empathy gap]}

\end{itemize}

\subsection{Data Sources for Tier 1 Calibration}

\begin{itemize}

\item \textbf{World Values Survey (WVS).} 
\url{https://www.worldvaluessurvey.org/}
\textit{[Trust, values, preferences across 100+ countries]}

\item \textbf{European Social Survey (ESS).} 
\url{https://www.europeansocialsurvey.org/}
\textit{[Social attitudes, wellbeing in Europe]}

\item \textbf{Gallup World Poll.} 
\url{https://www.gallup.com/analytics/318875/global-research.aspx}
\textit{[Life satisfaction, emotional wellbeing]}

\item \textbf{Panel Study of Income Dynamics (PSID).} 
\url{https://psidonline.isr.umich.edu/}
\textit{[Longitudinal income, health, wealth data]}

\item \textbf{Global Preference Survey (GPS).} 
Falk, A. et al. (2018). ``Global Evidence on Economic Preferences.'' 
\textit{Quarterly Journal of Economics}, 133(4), 1645-1692.
\textit{[Risk, time, social preferences across 76 countries]}

\end{itemize}

\subsection{Epistemology and Philosophy of Science}

\begin{itemize}

\item \textbf{Cartwright, N. (2007).} 
\textit{Hunting Causes and Using Them: Approaches in Philosophy and Economics}. 
Cambridge University Press.
\textit{[Causation, external validity]}

\item \textbf{Manski, C. F. (2013).} 
\textit{Public Policy in an Uncertain World}. 
Harvard University Press.
\textit{[Decision-making under partial knowledge]}

\item \textbf{Deaton, A. \& Cartwright, N. (2018).} 
``Understanding and Misunderstanding Randomized Controlled Trials.'' 
\textit{Social Science \& Medicine}, 210, 2-21.
\textit{[Limits of RCT evidence, mechanism importance]}

\end{itemize}


nd{document}
